\chapter{Complexity — Loop}
%%% generated by the blueprint pipeline from TCSlib.Complexity.TuringMachine.Build.Loop
%%%%%%%%%%%%%%%%%%%%%%%%%%%%%%%%%%%%%%%%%%%%%%%%%%%%%%%%%%%%%%%%%%%%%%%%%%%%%%%
\section{Overview}

This module provides the bounded-loop control centerpiece of the machine-construction
library for multi-tape Turing machines over the Boolean alphabet. Its public results are
the loop summation lemma, the configuration-level, decision, and result-bearing loop
combinators (a fuel machine lays down a binary round counter, and a body machine is
called once per round between visits to an anchor state), the emitting loop combinator
whose rounds append output chunks, and the install- and emit-mode clean-call bridges
that package any time-bounded function witness as a callable module with clean entry and
exit seams. The supporting material constructs and verifies the concrete loop
controllers phase by phase: fuel capture and counter setup, the stop-wrapped body calls,
the fixed-width counter countdown, payload replay, tracked-bank evaluation, native
interval cleanup, and result finalization.

%%%%%%%%%%%%%%%%%%%%%%%%%%%%%%%%%%%%%%%%%%%%%%%%%%%%%%%%%%%%%%%%%%%%%%%%%%%%%%%
\section{Declarations}

\begin{definition}[Standard word assignment for seam configurations]
\lean{Turing.stateWord}
\label{Turing.stateWord}
\leanok
For a tape count $k$ and a Boolean word $s$, the assignment of one word per work tape
that places $s$ on tape $0$ and the empty word on every other tape. It is the standard
tape content of a loop body's round-boundary (seam) configurations: round state on tape
zero, all scratch blank.
\end{definition}

\begin{theorem}[The loop summation lemma]
\lean{Turing.loop_run}
\label{Turing.loop_run}
\leanok
\uses{Turing.Cfg, Turing.MultiTapeTM, Turing.MultiTapeTM.runFrom, Turing.MultiTapeTM.runFrom_add}
\difficulty{5}
Let $M$ be a multi-tape Turing machine over the Boolean alphabet, let $c_0, \dots, c_N$
be configurations of $M$ on a common input, all with empty output, and let $a$ be a
Boolean predicate on round indices. Assume that from $c_N$ the machine halts within $B$
steps with output the single bit $\mathrm{false}$, and that for each $j < N$, within $B$
steps the machine either halts from $c_j$ with output the single bit $\mathrm{true}$
(when $a(j)$ holds) or reaches $c_{j+1}$ (when $a(j)$ fails). Then from $c_0$ the
machine halts within $(N+1) \cdot B$ steps, and its output is the single bit recording
whether $a(j)$ holds for some $j < N$.
\end{theorem}

\begin{lemma}[Liveness propagates to earlier times]
\lean{Turing.FinTM.loop_live_prefix}
\label{Turing.FinTM.loop_live_prefix}
\leanok
\uses{Turing.Cfg, Turing.MultiTapeTM, Turing.MultiTapeTM.runFrom, Turing.MultiTapeTM.runFrom_add, Turing.MultiTapeTM.runFrom_of_halt}
\difficulty{3}
Let $M$ be a multi-tape Turing machine over the Boolean alphabet and $c$ a configuration
of $M$. If after $t$ steps from $c$ the machine has not halted, then for every $u \le t$
the machine has not halted after $u$ steps either.
\end{lemma}

\begin{lemma}[Empty final output forces empty earlier outputs]
\lean{Turing.FinTM.loop_silent_prefix}
\label{Turing.FinTM.loop_silent_prefix}
\leanok
\uses{Turing.Cfg, Turing.MultiTapeTM, Turing.MultiTapeTM.output_prefix, Turing.MultiTapeTM.runFrom}
\difficulty{3}
Let $M$ be a multi-tape Turing machine over the Boolean alphabet and $c$ a configuration
of $M$. If the output after $t$ steps from $c$ is empty, then the output after $u$ steps
is empty for every $u \le t$, since output is append-only.
\end{lemma}

\begin{lemma}[First-halt refinement of a halting witness]
\lean{Turing.FinTM.loop_first_halt}
\label{Turing.FinTM.loop_first_halt}
\leanok
\uses{Turing.Cfg, Turing.Cfg.Halted, Turing.MultiTapeTM, Turing.MultiTapeTM.runFrom, Turing.MultiTapeTM.runFrom_add, Turing.MultiTapeTM.runFrom_of_halt, Turing.MultiTapeTM.runFrom_zero}
\difficulty{4}
Let $M$ be a multi-tape Turing machine over the Boolean alphabet and $c$ a configuration
whose control state is live. If $M$ has halted $t$ steps after $c$, then there is a time
$u$ with $0 < u \le t$ such that $M$ is not halted at any time strictly before $u$, is
halted at time $u$, and the configuration at time $u$ equals the configuration at time
$t$.
\end{lemma}

\begin{lemma}[Invariant holds along the whole orbit]
\lean{Turing.FinTM.loop_orbit_inv}
\label{Turing.FinTM.loop_orbit_inv}
\leanok
\difficulty{3}
Let $I$ be a binary predicate on Boolean words, let $F$ be a step function on Boolean
words indexed by an input word, and let $s_0$ assign an initial word to each input. If
$I(x, s_0(x))$ holds for every input $x$ and $I(x, s)$ implies $I(x, F(x, s))$, then
$I(x, F(x, \cdot)^{i}(s_0(x)))$ holds for every input $x$ and every iteration count $i$.
\end{lemma}

\begin{lemma}[Fuel output width is bounded by its running time]
\lean{Turing.FinTM.loop_fuel_width}
\label{Turing.FinTM.loop_fuel_width}
\leanok
\uses{Turing.FinTM, Turing.FinTM.ComputesFunInTime, Turing.MultiTapeTM.output_length_le}
\difficulty{2}
Let $F$ be a bundled finite machine over the Boolean alphabet that computes, in time
$T$, the function sending each input $x$ to the binary digit word of $R(|x|)$. Then for
every input $x$, the length of that digit word is at most $T(|x|)$, since a run can
append at most one output symbol per step.
\end{lemma}

\begin{lemma}[One input move advances position by at most one]
\lean{Turing.FinTM.loop_input_move_le}
\label{Turing.FinTM.loop_input_move_le}
\leanok
\uses{Turing.FinTM.moveInputPos_neg_val, Turing.moveInputPos, Turing.moveInputPos_pos_of_ne_right, Turing.moveInputPos_rightBoundary}
\difficulty{4}
For every clamped input-head position $p$ among the $n+2$ positions of an input of
length $n$ and every direction $m$ (left, stay, or right), the position reached by one
native input-head move in direction $m$ has value at most the value of $p$ plus one.
\end{lemma}

\begin{lemma}[Input displacement is bounded by elapsed time]
\lean{Turing.FinTM.loop_input_run_le}
\label{Turing.FinTM.loop_input_run_le}
\leanok
\uses{Turing.Action.apply, Turing.Cfg, Turing.Cfg.inputSymbol, Turing.Cfg.workTapeSymbols, Turing.FinTM.loop_input_move_le, Turing.MultiTapeTM, Turing.MultiTapeTM.runFrom, Turing.MultiTapeTM.runFrom_succ_eq_step', Turing.MultiTapeTM.step}
\difficulty{4}
Let $M$ be a multi-tape Turing machine over the Boolean alphabet and $c$ a configuration
of $M$. After running $t$ steps from $c$, the input-head position is at most the input-
head position of $c$ plus $t$.
\end{lemma}

\begin{lemma}[At most one output bit per step]
\lean{Turing.FinTM.loop_output_length_le}
\label{Turing.FinTM.loop_output_length_le}
\leanok
\uses{Turing.Cfg, Turing.FinTM.loopDebit, Turing.MultiTapeTM, Turing.MultiTapeTM.outputSymbol, Turing.MultiTapeTM.runFrom, Turing.MultiTapeTM.runFrom_succ_eq_step', Turing.MultiTapeTM.step, Turing.MultiTapeTM.step_output}
\difficulty{4}
Let $M$ be a multi-tape Turing machine over the Boolean alphabet and $c$ a configuration
of $M$. After running $t$ steps from $c$, the output length is at most the output length
of $c$ plus $t$.
\end{lemma}

\begin{lemma}[Bounded input rewind through a two-state scan]
\lean{Turing.FinTM.loop_rewind_bounded}
\label{Turing.FinTM.loop_rewind_bounded}
\leanok
\uses{Turing.Cfg, Turing.FinTM.controlAction, Turing.FinTM.controlAction_apply, Turing.FinTM.loopDebit, Turing.FinTM.moveInputPos_neg_val, Turing.FinTM.rewind_scan, Turing.MultiTapeTM, Turing.MultiTapeTM.runFrom, Turing.MultiTapeTM.runFrom_add, Turing.MultiTapeTM.step, Turing.moveInputPos}
\difficulty{5}
Let $M$ be a multi-tape Turing machine over the Boolean alphabet with two distinguished
control states: a start state whose every transition moves the input head left and
enters a scan state, and a scan state that keeps moving left while it reads an input
symbol and, on reading the left boundary blank, moves right and enters a given
destination state. Then from any configuration $c$ in the start state there is a time
$t$ at most the input-head position of $c$ plus two after which the machine sits in the
destination state with input head at position one and every other component of $c$
unchanged.
\end{lemma}

\begin{definition}[Fixed-width little-endian decrement]
\lean{Turing.FinTM.loopDebit}
\label{Turing.FinTM.loopDebit}
\leanok
For a Boolean word read as a little-endian binary counter, the pair of the decremented
word and a success flag: the lowest $\mathrm{true}$ cell is cleared and all lower cells
are set to $\mathrm{true}$, returning success; if no cell is $\mathrm{true}$
(underflow), every existing cell is set to $\mathrm{true}$ and the flag is
$\mathrm{false}$. The word length never changes and the empty word reports underflow.
\end{definition}

\begin{definition}[Length of the borrow scan]
\lean{Turing.FinTM.loopBorrowPos}
\label{Turing.FinTM.loopBorrowPos}
\leanok
For a Boolean word read as a little-endian counter, the number of low $\mathrm{false}$
cells traversed by a borrow; it equals the length of the maximal $\mathrm{false}$ prefix
of the word.
\end{definition}

\begin{lemma}[Borrow scan bounded by the width]
\lean{Turing.FinTM.loopBorrowPos_le}
\label{Turing.FinTM.loopBorrowPos_le}
\leanok
\uses{Turing.FinTM.loopBorrowPos}
\difficulty{3}
For every Boolean word $u$, the length of the maximal $\mathrm{false}$ prefix of $u$ is
at most the length of $u$.
\end{lemma}

\begin{lemma}[Decrement preserves the counter width]
\lean{Turing.FinTM.loopDebit_length}
\label{Turing.FinTM.loopDebit_length}
\leanok
\uses{Turing.FinTM.loopDebit}
\difficulty{3}
For every Boolean word $u$, the word produced by the fixed-width little-endian decrement
of $u$ has exactly the length of $u$, both on success and on underflow.
\end{lemma}

\begin{definition}[Value of a little-endian counter]
\lean{Turing.FinTM.loopValue}
\label{Turing.FinTM.loopValue}
\leanok
For a Boolean word $b_0 b_1 \cdots b_{m-1}$, the natural number $\sum_{i < m} b_i
2^{i}$, where each $\mathrm{true}$ cell contributes its power of two and high
$\mathrm{false}$ cells contribute nothing.
\end{definition}

\begin{lemma}[Binary digits evaluate to their number]
\lean{Turing.FinTM.loopValue_bits}
\label{Turing.FinTM.loopValue_bits}
\leanok
\uses{Turing.FinTM.loopValue}
\difficulty{4}
For every natural number $n$, the little-endian counter value of the standard binary
digit word of $n$ equals $n$.
\end{lemma}

\begin{lemma}[Decrement reduces the value by exactly one]
\lean{Turing.FinTM.loopDebit_value}
\label{Turing.FinTM.loopDebit_value}
\leanok
\uses{Turing.FinTM.loopDebit, Turing.FinTM.loopValue}
\difficulty{4}
For every Boolean word $u$ read as a little-endian counter: if the fixed-width decrement
of $u$ reports success then the value of the decremented word plus one equals the value
of $u$, and if it reports underflow then the value of $u$ is zero.
\end{lemma}

\begin{lemma}[Success exactly at positive values]
\lean{Turing.FinTM.loopDebit_success}
\label{Turing.FinTM.loopDebit_success}
\leanok
\uses{Turing.FinTM.loopDebit, Turing.FinTM.loopDebit_value, Turing.FinTM.loopValue}
\difficulty{3}
For every Boolean word $u$ read as a little-endian counter, the fixed-width decrement of
$u$ reports success if and only if the value of $u$ is positive.
\end{lemma}

\begin{lemma}[Iterated decrements preserve the width]
\lean{Turing.FinTM.loopDebit_iterate_length}
\label{Turing.FinTM.loopDebit_iterate_length}
\leanok
\uses{Turing.FinTM.loopDebit, Turing.FinTM.loopDebit_length}
\difficulty{3}
For every Boolean word $u$ and every iteration count $i$, applying the fixed-width
little-endian decrement $i$ times to $u$ yields a word of exactly the length of $u$.
\end{lemma}

\begin{lemma}[Counter value after i decrements]
\lean{Turing.FinTM.loopDebit_iterate_value}
\label{Turing.FinTM.loopDebit_iterate_value}
\leanok
\uses{Turing.FinTM.loopDebit, Turing.FinTM.loopDebit_success, Turing.FinTM.loopDebit_value, Turing.FinTM.loopValue, Turing.FinTM.loopValue_bits}
\difficulty{4}
For all natural numbers $R$ and $i$ with $i \le R$, applying the fixed-width little-
endian decrement $i$ times to the binary digit word of $R$ yields a word of counter
value exactly $R - i$.
\end{lemma}

\begin{lemma}[Reading the frontier cell of a buffer tape]
\lean{Turing.FinTM.loopBuffer_read}
\label{Turing.FinTM.loopBuffer_read}
\leanok
\uses{Turing.FinTM.bufferTape, Turing.FinTM.bufferTape_nat, Turing.FinTM.loopDebitTM}
\difficulty{2}
For all Boolean words $p$ and $b$, the buffer tape holding the concatenation of $p$ and
$b$ (cells $0, \dots$ hold the word, every other cell is blank) reads, at cell $|p|$,
the first bit of $b$, or blank when $b$ is empty.
\end{lemma}

\begin{lemma}[Overwriting one cell of a buffered word]
\lean{Turing.FinTM.loopBuffer_write}
\label{Turing.FinTM.loopBuffer_write}
\leanok
\uses{Turing.FinTM.bufferTape, Turing.FinTM.loopDebitTM}
\difficulty{4}
For all Boolean words $p$ and $b$ and bits $u$ and $v$: updating cell $|p|$ of the
buffer tape that holds $p$ followed by $u$ followed by $b$, writing the bit $v$ there,
produces exactly the buffer tape that holds $p$ followed by $v$ followed by $b$.
\end{lemma}

\begin{definition}[One-tape counter decrement machine]
\lean{Turing.FinTM.loopDebitTM}
\label{Turing.FinTM.loopDebitTM}
\leanok
\uses{Turing.FinTM, Turing.FinTM.controlAction}
A bundled finite machine over the Boolean alphabet with one work tape that performs a
fixed-width little-endian decrement followed by a rewind to cell zero. Its live control
states are a borrow state, rewind states remembering the success flag, and silent
waiting return states indexed by that flag; no transition emits physical output.
\end{definition}

\begin{definition}[Configuration of the decrement machine]
\lean{Turing.FinTM.loopDebitCfg}
\label{Turing.FinTM.loopDebitCfg}
\leanok
\uses{Turing.Cfg, Turing.FinTM.bufferTape, Turing.FinTM.loopDebitTM}
For an input word $x$, an input-head position $p$, a control state $q$, a work head
position $z$, and a counter word $u$, the configuration of the one-tape counter
decrement machine whose single work tape is the buffer tape of $u$, with head at $z$,
control $q$, input head at $p$, and empty output.
\end{definition}

\begin{lemma}[One borrow transition of the decrement machine]
\lean{Turing.FinTM.loopBorrow_step}
\label{Turing.FinTM.loopBorrow_step}
\leanok
\uses{Turing.Action.apply, Turing.Cfg.workTapeSymbols, Turing.FinTM.loopBuffer_read, Turing.FinTM.loopBuffer_write, Turing.FinTM.loopDebitCfg, Turing.FinTM.loopDebitTM, Turing.MultiTapeTM.step, Turing.moveInputPos_zero}
\difficulty{4}
Let $p$ and $b$ be Boolean words. One step of the one-tape counter decrement machine
from its borrow state, at cell $|p|$ of the buffer tape holding $p$ followed by $b$,
behaves as follows: on an empty $b$ it enters the rewind state with failure flag one
cell to the left without writing; on a leading $\mathrm{true}$ bit it writes
$\mathrm{false}$, moves left, and enters the rewind state with success flag; on a
leading $\mathrm{false}$ bit it writes $\mathrm{true}$, moves right, and stays in the
borrow state.
\end{lemma}

\begin{lemma}[The borrow phase in exact time]
\lean{Turing.FinTM.loopBorrow_run}
\label{Turing.FinTM.loopBorrow_run}
\leanok
\uses{Turing.FinTM.loopBorrowPos, Turing.FinTM.loopBorrow_step, Turing.FinTM.loopDebit, Turing.FinTM.loopDebitCfg, Turing.FinTM.loopDebitTM, Turing.MultiTapeTM.runFrom, Turing.MultiTapeTM.runFrom_succ_eq_step}
\difficulty{4}
Let $p$ and $u$ be Boolean words and let $j$ be the length of the maximal
$\mathrm{false}$ prefix of $u$. Starting the one-tape counter decrement machine in its
borrow state at cell $|p|$ of the buffer tape holding $p$ followed by $u$, after exactly
$j + 1$ steps it is in the rewind state carrying the decrement's success flag, at cell
$|p| + j - 1$, with tape contents $p$ followed by the decremented word.
\end{lemma}

\begin{lemma}[The rewind phase in exact time]
\lean{Turing.FinTM.loopBorrow_rewind}
\label{Turing.FinTM.loopBorrow_rewind}
\leanok
\uses{Turing.Action.apply, Turing.Cfg.workTapeSymbols, Turing.FinTM.bufferTape_left, Turing.FinTM.bufferTape_nat, Turing.FinTM.loopDebitCfg, Turing.FinTM.loopDebitTM, Turing.MultiTapeTM.runFrom, Turing.MultiTapeTM.runFrom_succ_eq_step, Turing.MultiTapeTM.runFrom_zero, Turing.MultiTapeTM.step, Turing.moveInputPos_zero}
\difficulty{4}
Let $u$ be a Boolean word, $b$ a flag bit, and $j \le |u|$. Starting the one-tape
counter decrement machine in its rewind state with flag $b$ at cell $j - 1$ over the
buffer tape of $u$, after exactly $j + 1$ steps it is in the return state indexed by
$b$, at cell zero, with the word $u$ unchanged.
\end{lemma}

\begin{lemma}[Complete decrement-and-rewind contract]
\lean{Turing.FinTM.loopBorrow_correct}
\label{Turing.FinTM.loopBorrow_correct}
\leanok
\uses{Turing.FinTM.loopBorrowPos, Turing.FinTM.loopBorrowPos_le, Turing.FinTM.loopBorrow_rewind, Turing.FinTM.loopBorrow_run, Turing.FinTM.loopDebit, Turing.FinTM.loopDebitCfg, Turing.FinTM.loopDebitTM, Turing.FinTM.loopDebit_length, Turing.MultiTapeTM.runFrom, Turing.MultiTapeTM.runFrom_add}
\difficulty{4}
Let $u$ be a Boolean word and $j$ the length of its maximal $\mathrm{false}$ prefix.
Then $2j + 2 \le 2|u| + 2$, and running the one-tape counter decrement machine for
exactly $2j + 2$ steps from its borrow state at cell zero over the buffer tape of $u$
ends in the return state indexed by the decrement's success flag, at cell zero, holding
the decremented word, with nothing emitted.
\end{lemma}

\begin{definition}[Stop wrapper for a loop body]
\lean{Turing.FinTM.loopBodyTM}
\label{Turing.FinTM.loopBodyTM}
\leanok
\uses{Turing.FinTM}
Given a bundled finite machine (the loop body) and a distinguished anchor control state,
the wrapper machine with one extra one-cell flag tape that executes the body but halts
silently at the next entry into the anchor, recording on the flag tape whether the stop
was an anchor return ($\mathrm{false}$) or a genuine body halt ($\mathrm{true}$). A set
release bit forces one body action even at the anchor, and every successor action clears
it; the body's output is retained in full.
\end{definition}

\begin{definition}[Embedding a body configuration into the stop wrapper]
\lean{Turing.FinTM.loopBodyCfg}
\label{Turing.FinTM.loopBodyCfg}
\leanok
\uses{Turing.Cfg, Turing.FinTM, Turing.FinTM.loopBodyTM}
For a configuration $c$ of the loop body, a release bit, and an optional flag bit, the
configuration of the stop wrapper whose body tapes carry the data of $c$, whose extra
flag tape stores the flag bit in its origin cell with head at the origin, and whose
control state is the control state of $c$ paired with the release bit.
\end{definition}

\begin{lemma}[Silent stop at an unreleased anchor]
\lean{Turing.FinTM.loopBody_stop}
\label{Turing.FinTM.loopBody_stop}
\leanok
\uses{Turing.Action.apply, Turing.Cfg, Turing.FinTM, Turing.FinTM.loopBodyCfg, Turing.FinTM.loopBodyTM, Turing.MultiTapeTM.step, Turing.moveInputPos_zero}
\difficulty{4}
Let $c$ be a configuration of the loop body whose control state is the anchor and whose
release bit is clear. One step of the stop wrapper from the embedding of $c$ halts the
wrapper, writes $\mathrm{false}$ on the flag cell, and changes no other component of the
embedded configuration.
\end{lemma}

\begin{lemma}[The wrapper executes one genuine body action]
\lean{Turing.FinTM.loopBody_step}
\label{Turing.FinTM.loopBody_step}
\leanok
\uses{Turing.Action, Turing.Action.apply, Turing.Cfg, Turing.Cfg.inputSymbol, Turing.Cfg.workTapeSymbols, Turing.FinTM, Turing.FinTM.loopBodyCfg, Turing.FinTM.loopBodyTM, Turing.MultiTapeTM.step}
\difficulty{5}
Let $c$ be a configuration of the loop body in a live control state $q$, and suppose it
is not the case that $q$ is the anchor with the release bit clear. Then one step of the
stop wrapper from the embedding of $c$ equals the embedding of the body's one-step
successor of $c$, with the release bit cleared and the flag set to $\mathrm{true}$
exactly when that body step is a halting transition.
\end{lemma}

\begin{lemma}[Exact simulation up to the first halt or anchor]
\lean{Turing.FinTM.loopBody_run}
\label{Turing.FinTM.loopBody_run}
\leanok
\uses{Turing.Cfg, Turing.FinTM, Turing.FinTM.loopBodyCfg, Turing.FinTM.loopBodyTM, Turing.FinTM.loopBody_step, Turing.MultiTapeTM.runFrom, Turing.MultiTapeTM.runFrom_succ_eq_step'}
\difficulty{4}
Let $c$ be a live configuration of the loop body, let $r$ be a release bit, and let $t$
be a time such that the body stays live before $t$ and, except possibly for a released
step at time zero, does not visit the anchor before $t$. Then the stop wrapper run for
$t$ steps from the embedding of $c$ with release bit $r$ equals the embedding of the
body run at time $t$, with the release bit consumed after the first action and the flag
set to $\mathrm{true}$ exactly when the body has halted at time $t$.
\end{lemma}

\begin{lemma}[Capture of the stopped body in an agreeing host]
\lean{Turing.FinTM.loopBody_capture}
\label{Turing.FinTM.loopBody_capture}
\leanok
\uses{Turing.Cfg, Turing.Cfg.Halted, Turing.FinTM, Turing.FinTM.LoopHostState, Turing.FinTM.loopBodyCfg, Turing.FinTM.loopBodyTM, Turing.FinTM.loopBody_run, Turing.MultiTapeTM, Turing.MultiTapeTM.runFrom, Turing.captureAction, Turing.captureCfg, Turing.capture_run}
\difficulty{4}
Let a host machine agree, on a family of embedded control states, with the capture
transformation of the stop wrapper of a loop body, and let $c$ be a live body
configuration together with a release bit and a time $t$ satisfying the liveness and no-
early-anchor guards. Then the host run for $t$ steps from the capture embedding of $c$
equals the capture embedding of the simulated stop-wrapper endpoint at time $t$,
including any output bit emitted by a halting transition.
\end{lemma}

\begin{definition}[Control states of the loop controller]
\lean{Turing.FinTM.LoopHostState}
\label{Turing.FinTM.LoopHostState}
\leanok
\uses{Turing.FinTM}
The finite control type of the loop host is a disjoint sum of the fuel machine's states,
the body-call states (a stopped-body state paired with a startup marker), and fourteen
controller phases.
\end{definition}

\begin{definition}[Fuel machine relocated past the body block]
\lean{Turing.FinTM.loopFuelSource}
\label{Turing.FinTM.loopFuelSource}
\leanok
\uses{Turing.FinTM, Turing.FinTM.rightAction, Turing.MultiTapeTM}
Given a loop body and a fuel machine $F$, the multi-tape machine with the combined tape
count whose transitions are those of $F$ acting on the last block of work tapes, while
the body tapes, the flag tape, and the counter tape are left untouched.
\end{definition}

\begin{definition}[Stopped body extended by counter and fuel tapes]
\lean{Turing.FinTM.loopBodySource}
\label{Turing.FinTM.loopBodySource}
\leanok
\uses{Turing.FinTM, Turing.FinTM.leftAction, Turing.FinTM.loopBodyTM, Turing.MultiTapeTM}
Given a loop body with anchor and a fuel machine, the multi-tape machine on the combined
tape count whose transitions are those of the body's stop wrapper acting on the leading
tapes, while the counter tape and the fuel residue tapes are inactive.
\end{definition}

\begin{definition}[Controller action on the three active tracks]
\lean{Turing.FinTM.loopControlAction}
\label{Turing.FinTM.loopControlAction}
\leanok
\uses{Turing.Action, Turing.FinTM, Turing.FinTM.LoopHostState}
For an input-head direction, an optional flag write, counter and payload write-and-move
pairs, an optional emitted bit, and a successor control state, the action of the loop
host that performs exactly these operations; it touches only the flag, counter, and
capture (payload) tapes, and every body and fuel tape is stationary under it.
\end{definition}

\begin{definition}[The concrete loop controller]
\lean{Turing.FinTM.loopHost}
\label{Turing.FinTM.loopHost}
\leanok
\uses{Turing.FinTM, Turing.FinTM.LoopHostState, Turing.FinTM.loopBodySource, Turing.FinTM.loopControlAction, Turing.FinTM.loopFuelSource, Turing.captureAction}
Given a loop body with anchor, a fuel machine, and an output mode, the bundled finite
machine that runs the fuel machine under capture, rewinds and copies the captured fuel
word into a fixed-width binary counter, rewinds the native input, and then repeatedly
calls the body between anchor visits: a startup call releases the first round free of
charge, each rejecting round decrements the counter in place, and the run ends either at
an accepting round -- emitting the verdict bit $\mathrm{true}$ in decision mode or
replaying the captured payload in find mode -- or at counter underflow, emitting
$\mathrm{false}$ in decision mode and nothing in find mode. Its fourteen controller
phases cover fuel rewind, counter copy, head rewinds, startup and round returns, borrow,
success and underflow rewinds, exhaustion, and payload replay.
\end{definition}

\begin{lemma}[Body calls agree with the capture discipline]
\lean{Turing.FinTM.loopHost_body_capture}
\label{Turing.FinTM.loopHost_body_capture}
\leanok
\uses{Turing.Cfg, Turing.Cfg.Halted, Turing.FinTM, Turing.FinTM.loopBodySource, Turing.FinTM.loopHost, Turing.MultiTapeTM.runFrom, Turing.captureCfg, Turing.capture_run}
\difficulty{1}
Let $c$ be a configuration of the extended stopped-body source on the loop host's tape
layout and $t$ a time before which that source does not halt. Then the loop host run for
$t$ steps from the capture embedding of $c$ (with the return phase determined by whether
the call is a startup call) equals the capture embedding of the source run at time $t$.
\end{lemma}

\begin{lemma}[Fuel phase captures all fuel emissions]
\lean{Turing.FinTM.loopHost_fuel_capture}
\label{Turing.FinTM.loopHost_fuel_capture}
\leanok
\uses{Turing.Cfg, Turing.Cfg.Halted, Turing.FinTM, Turing.FinTM.loopFuelSource, Turing.FinTM.loopHost, Turing.MultiTapeTM.runFrom, Turing.captureCfg, Turing.capture_run}
\difficulty{1}
Let $c$ be a configuration of the relocated fuel source on the loop host's tape layout
and $t$ a time before which that source does not halt. Then the loop host run for $t$
steps from the capture embedding of $c$ equals the capture embedding of the fuel source
run at time $t$.
\end{lemma}

\begin{lemma}[Initial configuration of the loop host]
\lean{Turing.FinTM.loopHost_init}
\label{Turing.FinTM.loopHost_init}
\leanok
\uses{Turing.Cfg.ofWords, Turing.FinTM, Turing.FinTM.loopFuelSource, Turing.FinTM.loopHost, Turing.MultiTapeTM.initCfg, Turing.captureCfg, Turing.initCfg_ofWords}
\difficulty{2}
For every input word $x$, the initial configuration of the loop host on $x$ equals the
capture embedding (with empty capture buffers and return phase zero) of the initial
configuration of the relocated fuel source on $x$.
\end{lemma}

\begin{lemma}[Idle controller actions are plain control actions]
\lean{Turing.FinTM.loopControl_idle}
\label{Turing.FinTM.loopControl_idle}
\leanok
\uses{Turing.FinTM, Turing.FinTM.LoopHostState, Turing.FinTM.controlAction, Turing.FinTM.loopControlAction}
\difficulty{1}
For every input direction and successor state, the loop-controller action that performs
no flag, counter, or payload operation and emits nothing equals the standard input-only
control action with that direction and successor.
\end{lemma}

\begin{lemma}[Phases four and five rewind the native input]
\lean{Turing.FinTM.loopHost_input_rewind}
\label{Turing.FinTM.loopHost_input_rewind}
\leanok
\uses{Turing.Cfg, Turing.FinTM, Turing.FinTM.loopControl_idle, Turing.FinTM.loopHost, Turing.FinTM.loop_rewind_bounded, Turing.MultiTapeTM.runFrom}
\difficulty{3}
Let $c$ be a configuration of the loop host in controller phase four. Then there is a
time $t$ at most the input-head position of $c$ plus two after which the host sits in
the startup body-call state with input head at position one and every other component of
$c$ unchanged.
\end{lemma}

\begin{definition}[Controller frame with explicit active tracks]
\lean{Turing.FinTM.loopFrame}
\label{Turing.FinTM.loopFrame}
\leanok
\uses{Turing.Cfg, Turing.FinTM, Turing.FinTM.LoopHostState}
For a base configuration of the loop host, a control state, an input position, three
tape functions for the flag, counter, and payload tracks, two head positions, and an
output word, the configuration whose three active controller tracks are given by these
data and whose remaining tapes and heads are those of the base configuration.
\end{definition}

\begin{definition}[Optional single-cell tape write]
\lean{Turing.FinTM.loopWrite}
\label{Turing.FinTM.loopWrite}
\leanok
For a tape, a head position, and an optional write instruction, the updated tape: with
no instruction the tape is unchanged, and with a write instruction exactly the cell
under the head is replaced by the written symbol.
\end{definition}

\begin{lemma}[Effect of a controller action on a frame]
\lean{Turing.FinTM.loopControl_apply}
\label{Turing.FinTM.loopControl_apply}
\leanok
\uses{Turing.Action.apply, Turing.Cfg, Turing.FinTM, Turing.FinTM.LoopHostState, Turing.FinTM.loopControlAction, Turing.FinTM.loopFrame, Turing.FinTM.loopReplayTM, Turing.FinTM.loopWrite, Turing.moveInputPos}
\difficulty{4}
Let a controller frame of the loop host be given, together with a controller action
specified by an input direction, an optional flag write, counter and payload write-and-
move pairs, an optional emitted bit, and a successor state. Applying the action to the
frame yields the frame with the successor state, the moved input head, the three tracks
updated by their optional writes at their current heads, the counter and payload heads
moved as instructed, and the emitted bit appended to the output.
\end{lemma}

\begin{definition}[One-tape payload replay machine]
\lean{Turing.FinTM.loopReplayTM}
\label{Turing.FinTM.loopReplayTM}
\leanok
\uses{Turing.FinTM}
A bundled finite machine over the Boolean alphabet with one work tape and a single
control state that, while its work cell holds a bit, emits that bit to the physical
output and moves right, and halts silently upon reading the right blank.
\end{definition}

\begin{definition}[Configuration of the replay machine]
\lean{Turing.FinTM.loopReplayCfg}
\label{Turing.FinTM.loopReplayCfg}
\leanok
\uses{Turing.Cfg, Turing.FinTM.bufferTape}
For an input word, an input-head position $p$, an optional control state, a head
position $z$, a stored word $w$, and an output prefix, the configuration of the one-tape
replay machine whose work tape is the buffer tape of $w$ with head at $z$, input head at
$p$, and the given output.
\end{definition}

\begin{lemma}[One replay step emits the current bit]
\lean{Turing.FinTM.loopReplay_step}
\label{Turing.FinTM.loopReplay_step}
\leanok
\uses{Turing.Action.apply, Turing.Cfg.workTapeSymbols, Turing.FinTM.loopBuffer_read, Turing.FinTM.loopReplayCfg, Turing.FinTM.loopReplayTM, Turing.MultiTapeTM.step, Turing.moveInputPos_zero}
\difficulty{4}
Let $p$ and $r$ be Boolean words. One step of the one-tape replay machine from its live
state at cell $|p|$ of the buffer tape holding $p$ followed by $r$ either halts without
writing when $r$ is empty, or, when $r$ starts with a bit $b$, appends $b$ to the output
and advances the head by one, leaving the stored word unchanged.
\end{lemma}

\begin{lemma}[Replay emits the remaining suffix exactly]
\lean{Turing.FinTM.loopReplay_run}
\label{Turing.FinTM.loopReplay_run}
\leanok
\uses{Turing.FinTM.loopReplayCfg, Turing.FinTM.loopReplayTM, Turing.FinTM.loopReplay_step, Turing.MultiTapeTM.runFrom, Turing.MultiTapeTM.runFrom_succ_eq_step}
\difficulty{4}
Let $p$ and $r$ be Boolean words and let an output prefix be given. After running the
one-tape replay machine for exactly $|r| + 1$ steps from its live state at cell $|p|$ of
the buffer tape holding $p$ followed by $r$, the machine is halted with the head at cell
$|p| + |r|$, the stored word is unchanged, and the output is extended by exactly $r$ in
order.
\end{lemma}

\begin{lemma}[Payload-only actions are right-block extensions]
\lean{Turing.FinTM.loopControl_payload}
\label{Turing.FinTM.loopControl_payload}
\leanok
\uses{Turing.Action, Turing.FinTM, Turing.FinTM.LoopHostState, Turing.FinTM.loopControlAction, Turing.FinTM.rightAction}
\difficulty{4}
For every head direction $d$, optional emitted bit, and successor state, the loop-
controller action that only moves the payload head by $d$ and emits that bit equals the
right-block extension (leaving all earlier tapes stationary) of the corresponding one-
tape action.
\end{lemma}

\begin{lemma}[Phase thirteen replays the captured payload]
\lean{Turing.FinTM.loopHost_replay}
\label{Turing.FinTM.loopHost_replay}
\leanok
\uses{Turing.FinTM, Turing.FinTM.loopControlAction, Turing.FinTM.loopControl_payload, Turing.FinTM.loopHost, Turing.FinTM.loopReplayCfg, Turing.FinTM.loopReplayTM, Turing.FinTM.loopReplay_run, Turing.FinTM.rightAction, Turing.FinTM.rightCfg, Turing.FinTM.rightCfg_run, Turing.MultiTapeTM.runFrom}
\difficulty{5}
Let $w$ be a captured word, let an output prefix and arbitrary completed body and fuel
tapes with their head positions be given, and consider the loop host in its replay phase
with the payload head at cell zero of the buffer tape of $w$. After running the host for
exactly $|w| + 1$ steps, the host is halted, its output is extended by exactly $w$, and
the stored word, all inactive tapes, and all inactive heads are unchanged.
\end{lemma}

\begin{definition}[Fuel configuration on its relocated block]
\lean{Turing.FinTM.loopFuelCfg}
\label{Turing.FinTM.loopFuelCfg}
\leanok
\uses{Turing.Cfg, Turing.FinTM, Turing.FinTM.rightCfg}
For a configuration $c$ of the fuel machine, the configuration of the relocated fuel
source that carries the data of $c$ on the fuel block while the body tapes, the flag
tape, and the counter tape are blank with heads at the origin.
\end{definition}

\begin{lemma}[Relocation preserves fuel runs]
\lean{Turing.FinTM.loopFuel_run}
\label{Turing.FinTM.loopFuel_run}
\leanok
\uses{Turing.Cfg, Turing.FinTM, Turing.FinTM.loopFuelCfg, Turing.FinTM.loopFuelSource, Turing.FinTM.rightAction, Turing.FinTM.rightCfg_run, Turing.MultiTapeTM, Turing.MultiTapeTM.runFrom}
\difficulty{4}
Let $c$ be a configuration of the fuel machine and $t$ a time. Running the relocated
fuel source for $t$ steps from the relocated embedding of $c$ equals the relocated
embedding of the fuel machine's run from $c$ at time $t$.
\end{lemma}

\begin{lemma}[Relocated fuel starts blank]
\lean{Turing.FinTM.loopFuel_init}
\label{Turing.FinTM.loopFuel_init}
\leanok
\uses{Turing.Cfg.init, Turing.FinTM, Turing.FinTM.loopFuelCfg, Turing.FinTM.loopFuelSource, Turing.FinTM.rightCfg, Turing.MultiTapeTM.initCfg}
\difficulty{4}
For every input word $x$, the initial configuration of the relocated fuel source on $x$
equals the relocated embedding of the fuel machine's initial configuration on $x$.
\end{lemma}

\begin{lemma}[The payload track reads its parameterized tape]
\lean{Turing.FinTM.loopFrame_payload}
\label{Turing.FinTM.loopFrame_payload}
\leanok
\uses{Turing.Cfg, Turing.Cfg.workTapeSymbols, Turing.FinTM, Turing.FinTM.LoopHostState, Turing.FinTM.loopFrame}
\difficulty{2}
For every controller frame of the loop host with payload tape $f$ and payload head $h$,
the symbol scanned on the capture (payload) track of the frame is exactly $f(h)$.
\end{lemma}

\begin{lemma}[The counter track reads its parameterized tape]
\lean{Turing.FinTM.loopFrame_counter}
\label{Turing.FinTM.loopFrame_counter}
\leanok
\uses{Turing.Cfg, Turing.Cfg.workTapeSymbols, Turing.FinTM, Turing.FinTM.LoopHostState, Turing.FinTM.loopFrame}
\difficulty{1}
For every controller frame of the loop host with counter tape $f$ and counter head $h$,
the symbol scanned on the counter track of the frame is exactly $f(h)$.
\end{lemma}

\begin{lemma}[Phase one rewinds the capture head]
\lean{Turing.FinTM.loopHost_fuel_rewind}
\label{Turing.FinTM.loopHost_fuel_rewind}
\leanok
\uses{Turing.Action.apply, Turing.Cfg, Turing.Cfg.workTapeSymbols, Turing.FinTM, Turing.FinTM.LoopHostState, Turing.FinTM.bufferTape, Turing.FinTM.bufferTape_left, Turing.FinTM.bufferTape_nat, Turing.FinTM.loopControlAction, Turing.FinTM.loopControl_apply, Turing.FinTM.loopCopyTape, Turing.FinTM.loopFrame, Turing.FinTM.loopFrame_payload, Turing.FinTM.loopHost, Turing.FinTM.loopWrite, Turing.MultiTapeTM.runFrom, Turing.MultiTapeTM.runFrom_succ_eq_step, Turing.MultiTapeTM.runFrom_zero, Turing.MultiTapeTM.step}
\difficulty{5}
Let $w$ be a stored word and $j \le |w|$. Starting the loop host in its fuel-rewind
phase on a controller frame whose payload track is the buffer tape of $w$ with payload
head at cell $j - 1$, after exactly $j + 1$ steps the host is in the counter-copy phase
with the payload head at cell zero and every other frame component unchanged.
\end{lemma}

\begin{definition}[Capture tape during fuel copying]
\lean{Turing.FinTM.loopCopyTape}
\label{Turing.FinTM.loopCopyTape}
\leanok
\uses{Turing.FinTM.bufferTape}
For Boolean words $p$ (the processed prefix) and $r$ (the remaining suffix), the tape
that is blank on the first $|p|$ cells and agrees with the buffer tape of the
concatenation of $p$ and $r$ everywhere else.
\end{definition}

\begin{lemma}[The copying frontier reads the next fuel bit]
\lean{Turing.FinTM.loopCopy_read}
\label{Turing.FinTM.loopCopy_read}
\leanok
\uses{Turing.FinTM.loopBuffer_read, Turing.FinTM.loopCopyTape}
\difficulty{1}
For all Boolean words $p$ and $r$, the partially cleared copying tape with processed
prefix $p$ and suffix $r$ reads, at cell $|p|$, the first bit of $r$, or blank when $r$
is empty.
\end{lemma}

\begin{lemma}[Clearing one copied cell extends the prefix]
\lean{Turing.FinTM.loopCopy_erase}
\label{Turing.FinTM.loopCopy_erase}
\leanok
\uses{Turing.FinTM.loopCopyTape}
\difficulty{3}
For all Boolean words $p$ and $r$ and every bit $b$, blanking cell $|p|$ of the
partially cleared copying tape with processed prefix $p$ and suffix $b$ followed by $r$
yields the partially cleared copying tape with processed prefix $p$ extended by $b$ and
suffix $r$.
\end{lemma}

\begin{lemma}[Copying starts from the full fuel buffer]
\lean{Turing.FinTM.loopCopy_initial}
\label{Turing.FinTM.loopCopy_initial}
\leanok
\uses{Turing.FinTM.bufferTape, Turing.FinTM.loopCopyTape}
\difficulty{2}
For every Boolean word $w$, the partially cleared copying tape with empty processed
prefix and suffix $w$ equals the buffer tape of $w$.
\end{lemma}

\begin{lemma}[Copying ends with a blank capture tape]
\lean{Turing.FinTM.loopCopy_final}
\label{Turing.FinTM.loopCopy_final}
\leanok
\uses{Turing.FinTM.bufferTape, Turing.FinTM.loopCopyTape}
\difficulty{2}
For every Boolean word $w$, the partially cleared copying tape with processed prefix $w$
and empty suffix equals the blank buffer tape.
\end{lemma}

\begin{lemma}[Phase two copies fuel into the counter]
\lean{Turing.FinTM.loopHost_fuel_copy}
\label{Turing.FinTM.loopHost_fuel_copy}
\leanok
\uses{Turing.Action.apply, Turing.Cfg, Turing.Cfg.workTapeSymbols, Turing.FinTM, Turing.FinTM.LoopHostState, Turing.FinTM.bufferTape, Turing.FinTM.bufferTape_append, Turing.FinTM.loopControlAction, Turing.FinTM.loopControl_apply, Turing.FinTM.loopCopyTape, Turing.FinTM.loopCopy_erase, Turing.FinTM.loopCopy_final, Turing.FinTM.loopCopy_read, Turing.FinTM.loopFrame, Turing.FinTM.loopFrame_payload, Turing.FinTM.loopHost, Turing.FinTM.loopWrite, Turing.MultiTapeTM.runFrom, Turing.MultiTapeTM.runFrom_succ_eq_step, Turing.MultiTapeTM.runFrom_zero, Turing.MultiTapeTM.step}
\difficulty{5}
Let $p$ be the already-copied prefix and $r$ the remaining fuel suffix. Starting the
loop host in its counter-copy phase on a controller frame whose counter track is the
buffer tape of $p$ and whose payload track is the partially cleared copying tape, with
both heads at cell $|p|$, after exactly $|r| + 1$ steps the host is in the synchronized-
rewind phase with the counter track holding the full fuel word, the payload track blank,
and both heads one cell left of the word's right end.
\end{lemma}

\begin{lemma}[Phase three rewinds counter and capture together]
\lean{Turing.FinTM.loopHost_fuel_return}
\label{Turing.FinTM.loopHost_fuel_return}
\leanok
\uses{Turing.Action.apply, Turing.Cfg, Turing.Cfg.workTapeSymbols, Turing.FinTM, Turing.FinTM.LoopHostState, Turing.FinTM.bufferTape, Turing.FinTM.bufferTape_left, Turing.FinTM.bufferTape_nat, Turing.FinTM.loopControlAction, Turing.FinTM.loopControl_apply, Turing.FinTM.loopFrame, Turing.FinTM.loopFrame_counter, Turing.FinTM.loopHost, Turing.FinTM.loopWrite, Turing.MultiTapeTM.runFrom, Turing.MultiTapeTM.runFrom_succ_eq_step, Turing.MultiTapeTM.runFrom_zero, Turing.MultiTapeTM.step}
\difficulty{5}
Let $w$ be the copied counter word and $j \le |w|$. Starting the loop host in its
synchronized-rewind phase on a controller frame whose counter track is the buffer tape
of $w$ and whose payload track is blank, with both heads at cell $j - 1$, after exactly
$j + 1$ steps the host is in the input-rewind phase with both heads at cell zero and
every other frame component unchanged.
\end{lemma}

\begin{lemma}[Fuel setup in exactly 3w+4 steps]
\lean{Turing.FinTM.loopHost_fuel_setup}
\label{Turing.FinTM.loopHost_fuel_setup}
\leanok
\uses{Turing.Action.apply, Turing.Cfg, Turing.FinTM, Turing.FinTM.LoopHostState, Turing.FinTM.bufferTape, Turing.FinTM.loopControlAction, Turing.FinTM.loopControl_apply, Turing.FinTM.loopCopy_initial, Turing.FinTM.loopFrame, Turing.FinTM.loopHost, Turing.FinTM.loopHost_fuel_copy, Turing.FinTM.loopHost_fuel_return, Turing.FinTM.loopHost_fuel_rewind, Turing.FinTM.loopWrite, Turing.MultiTapeTM.runFrom, Turing.MultiTapeTM.runFrom_add, Turing.MultiTapeTM.runFrom_succ_eq_step, Turing.MultiTapeTM.step}
\difficulty{4}
Let $w$ be the captured fuel word. Starting the loop host in controller phase zero on a
frame whose counter track is blank and whose payload track is the buffer tape of $w$
with payload head at cell $|w|$, after exactly $3|w| + 4$ steps the host is in the
input-rewind phase with the counter track holding $w$, the payload track blank, both
heads at cell zero, and all inactive data unchanged. The empty word is included.
\end{lemma}

\begin{definition}[Captured fuel endpoint of the host]
\lean{Turing.FinTM.loopFuelCaptured}
\label{Turing.FinTM.loopFuelCaptured}
\leanok
\uses{Turing.Cfg, Turing.FinTM, Turing.FinTM.LoopHostState, Turing.FinTM.loopFuelCfg, Turing.captureCfg}
For a configuration $c$ of the fuel machine, the loop-host configuration obtained by
capture-embedding (with empty buffers and return phase zero) the relocated embedding of
$c$; it retains the complete fuel work residue.
\end{definition}

\begin{definition}[Prepared startup configuration]
\lean{Turing.FinTM.loopReady}
\label{Turing.FinTM.loopReady}
\leanok
\uses{Turing.Cfg, Turing.FinTM, Turing.FinTM.LoopHostState, Turing.FinTM.bufferTape, Turing.FinTM.loopFrame, Turing.FinTM.loopFuelCaptured}
For a halted fuel configuration $c$, the loop-host configuration whose control is the
startup body-call state, whose input head is at one, whose flag and payload tracks are
blank, whose counter track holds the output word of $c$ with both active heads at zero,
whose output is empty, and which retains the completed fuel residue.
\end{definition}

\begin{lemma}[The captured fuel endpoint as a controller frame]
\lean{Turing.FinTM.loopFuelCaptured_frame}
\label{Turing.FinTM.loopFuelCaptured_frame}
\leanok
\uses{Turing.Cfg, Turing.FinTM, Turing.FinTM.bufferTape, Turing.FinTM.loopFrame, Turing.FinTM.loopFuelCaptured, Turing.FinTM.loopFuelCfg, Turing.FinTM.rightCfg, Turing.captureCfg}
\difficulty{5}
Let $c$ be a halted configuration of the fuel machine. The captured fuel endpoint of the
loop host equals the controller frame over itself in phase zero whose flag and counter
tracks are blank, whose payload track is the buffer tape of the output of $c$ with
payload head at its length, with the input head of $c$ and empty output.
\end{lemma}

\begin{lemma}[Preparation reaches startup within 5T+7 steps]
\lean{Turing.FinTM.loopHost_prepare}
\label{Turing.FinTM.loopHost_prepare}
\leanok
\uses{Turing.Cfg, Turing.Cfg.Halted, Turing.Cfg.init, Turing.FinTM, Turing.FinTM.ComputesFunInTime, Turing.FinTM.bufferTape, Turing.FinTM.loopFrame, Turing.FinTM.loopFuelCaptured, Turing.FinTM.loopFuelCaptured_frame, Turing.FinTM.loopFuelCfg, Turing.FinTM.loopFuel_init, Turing.FinTM.loopFuel_run, Turing.FinTM.loopHost, Turing.FinTM.loopHost_fuel_capture, Turing.FinTM.loopHost_fuel_setup, Turing.FinTM.loopHost_init, Turing.FinTM.loopHost_input_rewind, Turing.FinTM.loopReady, Turing.FinTM.loop_first_halt, Turing.FinTM.loop_fuel_width, Turing.FinTM.loop_input_run_le, Turing.FinTM.rightCfg, Turing.MultiTapeTM.initCfg, Turing.MultiTapeTM.runFrom, Turing.MultiTapeTM.runFrom_add}
\difficulty{5}
Let the fuel machine compute, in time $T$, the function sending each input $x$ to the
binary digit word of $R(|x|)$. Then for every input $x$ there are a halted fuel
configuration $c$ with output exactly that digit word and a time $t \le 5\,T(|x|) + 7$
such that the loop host, run for $t$ steps from its initial configuration on $x$,
reaches the prepared startup configuration built from $c$.
\end{lemma}

\begin{definition}[Stopped body padded by counter and fuel residue]
\lean{Turing.FinTM.loopBodyPadded}
\label{Turing.FinTM.loopBodyPadded}
\leanok
\uses{Turing.Cfg, Turing.FinTM, Turing.FinTM.bufferTape, Turing.FinTM.leftCfg, Turing.FinTM.loopBodyCfg}
For a body configuration $c$, a release bit, an optional flag bit, a counter word $w$,
and a fuel configuration, the configuration of the extended stopped-body source whose
leading tapes embed $c$ with its flag data, whose counter tape is the buffer tape of $w$
with head at the origin, and whose fuel block carries the fuel configuration's tapes and
heads.
\end{definition}

\begin{definition}[A body call inside the capturing host]
\lean{Turing.FinTM.loopCall}
\label{Turing.FinTM.loopCall}
\leanok
\uses{Turing.Cfg, Turing.FinTM, Turing.FinTM.LoopHostState, Turing.FinTM.loopBodyPadded, Turing.captureCfg}
For a startup marker, a body configuration $c$, a release bit, an optional flag bit, a
counter word, and a fuel configuration, the loop-host configuration obtained by capture-
embedding the padded stopped-body configuration, with the return phase chosen according
to the startup marker.
\end{definition}

\begin{lemma}[Padded source simulates the stop wrapper]
\lean{Turing.FinTM.loopBodySource_run}
\label{Turing.FinTM.loopBodySource_run}
\leanok
\uses{Turing.Cfg, Turing.FinTM, Turing.FinTM.leftCfg, Turing.FinTM.leftCfg_run, Turing.FinTM.loopBodySource, Turing.FinTM.loopBodyTM, Turing.MultiTapeTM.runFrom}
\difficulty{1}
Let $c$ be a configuration of the body's stop wrapper, and let arbitrary counter and
fuel tapes and heads be given. Running the extended stopped-body source for $t$ steps
from the left-block embedding of $c$ over those inactive tapes equals the left-block
embedding of the stop wrapper's run at time $t$ over the same inactive tapes.
\end{lemma}

\begin{lemma}[Anchor returns are captured in one extra step]
\lean{Turing.FinTM.loopHost_anchor_return}
\label{Turing.FinTM.loopHost_anchor_return}
\leanok
\uses{Turing.Cfg, Turing.Cfg.Halted, Turing.FinTM, Turing.FinTM.leftCfg, Turing.FinTM.loopBodyCfg, Turing.FinTM.loopBodyPadded, Turing.FinTM.loopBodySource_run, Turing.FinTM.loopBodyTM, Turing.FinTM.loopBody_run, Turing.FinTM.loopBody_stop, Turing.FinTM.loopCall, Turing.FinTM.loopHost, Turing.FinTM.loopHost_body_capture, Turing.FinTM.loop_live_prefix, Turing.MultiTapeTM.runFrom, Turing.MultiTapeTM.runFrom_succ_eq_step'}
\difficulty{5}
Let $c$ be a body configuration, $t$ a time at which the body run from $c$ sits in the
anchor state, with no earlier anchor visit except possibly a released step at time zero,
and with the release bit clear whenever $t = 0$. Then the loop host, run for $t + 1$
steps from the body call built on $c$, reaches the stopped call whose body data is the
body run at time $t$ with control removed, release clear, and stop flag
$\mathrm{false}$, with the counter word and fuel residue unchanged.
\end{lemma}

\begin{lemma}[Prepared startup is the canonical startup call]
\lean{Turing.FinTM.loopReady_call}
\label{Turing.FinTM.loopReady_call}
\leanok
\uses{Turing.Cfg, Turing.Cfg.init, Turing.FinTM, Turing.FinTM.leftCfg, Turing.FinTM.loopBodyCfg, Turing.FinTM.loopBodyPadded, Turing.FinTM.loopCall, Turing.FinTM.loopFrame, Turing.FinTM.loopFuelCaptured, Turing.FinTM.loopFuelCfg, Turing.FinTM.loopReady, Turing.FinTM.rightCfg, Turing.MultiTapeTM.initCfg, Turing.captureCfg}
\difficulty{4}
For every halted fuel configuration $f$, the prepared startup configuration of the loop
host equals the startup body call built on the body's initial configuration with release
clear, no flag, counter word the output of $f$, and fuel residue $f$.
\end{lemma}

\begin{lemma}[Phase six releases the first anchor]
\lean{Turing.FinTM.loopHost_release}
\label{Turing.FinTM.loopHost_release}
\leanok
\uses{Turing.Action.apply, Turing.Cfg, Turing.FinTM, Turing.FinTM.leftCfg, Turing.FinTM.loopBodyCfg, Turing.FinTM.loopBodyPadded, Turing.FinTM.loopCall, Turing.FinTM.loopControlAction, Turing.FinTM.loopHost, Turing.MultiTapeTM.step, Turing.captureCfg, Turing.moveInputPos_zero}
\difficulty{4}
Let $c$ be a body configuration, $w$ a counter word, and $f$ a fuel configuration. One
step of the loop host from the startup call on $c$ with control removed, release clear,
and flag $\mathrm{false}$ yields the active (non-startup) call on $c$ with control the
anchor, release set, and no flag, changing no body, counter, or fuel data.
\end{lemma}

\begin{lemma}[Startup reaches the first released candidate]
\lean{Turing.FinTM.loopHost_start}
\label{Turing.FinTM.loopHost_start}
\leanok
\uses{Turing.Cfg, Turing.Cfg.ofWords, Turing.FinTM, Turing.FinTM.loopCall, Turing.FinTM.loopHost, Turing.FinTM.loopHost_anchor_return, Turing.FinTM.loopHost_release, Turing.FinTM.loopReady, Turing.FinTM.loopReady_call, Turing.MultiTapeTM.initCfg, Turing.MultiTapeTM.runFrom, Turing.MultiTapeTM.runFrom_succ_eq_step', Turing.stateWord}
\difficulty{3}
Let $s$ be a word, $f$ a fuel configuration, and $t$ a time such that the body, run from
its initial configuration, first visits the anchor at time $t$, arriving at the clean
seam that carries $s$ on work tape zero with all other work tapes blank. Then the loop
host, run for $t + 2$ steps from the prepared startup configuration, reaches the active
released call on that seam with counter word the output of $f$ and fuel residue $f$.
\end{lemma}

\begin{lemma}[A genuine halt returns with the true flag]
\lean{Turing.FinTM.loopHost_halt_return}
\label{Turing.FinTM.loopHost_halt_return}
\leanok
\uses{Turing.Cfg, Turing.Cfg.Halted, Turing.FinTM, Turing.FinTM.leftCfg, Turing.FinTM.loopBodyCfg, Turing.FinTM.loopBodyPadded, Turing.FinTM.loopBodySource_run, Turing.FinTM.loopBody_run, Turing.FinTM.loopCall, Turing.FinTM.loopHost, Turing.FinTM.loopHost_body_capture, Turing.MultiTapeTM.runFrom}
\difficulty{4}
Let $c$ be a body configuration and $t > 0$ a time such that the body stays live before
$t$, does not visit the anchor strictly between times zero and $t$, and is halted at
time $t$. Then the loop host, run for $t$ steps from the active released call on $c$,
reaches the stopped call on the body run at time $t$ with release clear and stop flag
$\mathrm{true}$, the complete body output captured, including any bit emitted by the
halting transition.
\end{lemma}

\begin{lemma}[A body call as a controller frame]
\lean{Turing.FinTM.loopCall_frame}
\label{Turing.FinTM.loopCall_frame}
\leanok
\uses{Turing.Cfg, Turing.FinTM, Turing.FinTM.bufferTape, Turing.FinTM.leftCfg, Turing.FinTM.loopBodyCfg, Turing.FinTM.loopBodyPadded, Turing.FinTM.loopCall, Turing.FinTM.loopFrame, Turing.captureCfg}
\difficulty{4}
For every body call of the loop host, built on a body configuration $c$ with flag bit,
counter word $w$, and fuel residue, the call equals the controller frame over itself
whose flag track stores the flag bit at the origin, whose counter track is the buffer
tape of $w$ with head at zero, and whose payload track is the buffer tape of the output
of $c$ with head at that output's length, with empty physical output.
\end{lemma}

\begin{lemma}[Reframing changes control, flag, and counter only]
\lean{Turing.FinTM.loopCall_reframe}
\label{Turing.FinTM.loopCall_reframe}
\leanok
\uses{Turing.Cfg, Turing.FinTM, Turing.FinTM.bufferTape, Turing.FinTM.leftCfg, Turing.FinTM.loopBodyCfg, Turing.FinTM.loopBodyPadded, Turing.FinTM.loopCall, Turing.FinTM.loopFrame, Turing.captureCfg}
\difficulty{4}
Let $c$ be a body configuration with counter word $w$ and fuel residue, and let new
startup, release, flag, counter-word, and control-state parameters be given. The
controller frame built over the original call with the new call's control state, the new
flag at the origin, and the buffer tape of the new counter word equals the call built on
$c$ with the replaced control state and the new parameters.
\end{lemma}

\begin{lemma}[One borrow step in the actual host]
\lean{Turing.FinTM.loopHost_borrow_step}
\label{Turing.FinTM.loopHost_borrow_step}
\leanok
\uses{Turing.Action.apply, Turing.Cfg, Turing.Cfg.workTapeSymbols, Turing.FinTM, Turing.FinTM.LoopHostState, Turing.FinTM.bufferTape, Turing.FinTM.loopBuffer_read, Turing.FinTM.loopBuffer_write, Turing.FinTM.loopControlAction, Turing.FinTM.loopControl_apply, Turing.FinTM.loopFrame, Turing.FinTM.loopFrame_counter, Turing.FinTM.loopHost, Turing.FinTM.loopWrite, Turing.MultiTapeTM.step}
\difficulty{4}
Let $p$ and $r$ be Boolean words. One step of the loop host from its borrow phase on a
frame whose counter track is the buffer tape of $p$ followed by $r$ with counter head at
$|p|$ behaves as the standalone decrement machine: an empty $r$ enters the underflow-
rewind phase moving left; a leading $\mathrm{true}$ writes $\mathrm{false}$ and enters
the success-rewind phase moving left; a leading $\mathrm{false}$ writes $\mathrm{true}$
and continues the borrow moving right. Only the counter track changes.
\end{lemma}

\begin{lemma}[The borrow scan in the actual host]
\lean{Turing.FinTM.loopHost_borrow_run}
\label{Turing.FinTM.loopHost_borrow_run}
\leanok
\uses{Turing.Cfg, Turing.FinTM, Turing.FinTM.LoopHostState, Turing.FinTM.bufferTape, Turing.FinTM.loopBorrowPos, Turing.FinTM.loopDebit, Turing.FinTM.loopFrame, Turing.FinTM.loopHost, Turing.FinTM.loopHost_borrow_step, Turing.MultiTapeTM.runFrom, Turing.MultiTapeTM.runFrom_succ_eq_step}
\difficulty{4}
Let $p$ and $w$ be Boolean words and $j$ the length of the maximal $\mathrm{false}$
prefix of $w$. Running the loop host for exactly $j + 1$ steps from its borrow phase on
a frame whose counter track holds $p$ followed by $w$ with head at $|p|$ reaches the
success- or underflow-rewind phase according to the decrement's flag, with the counter
track holding $p$ followed by the decremented word and head at $|p| + j - 1$; all non-
counter tracks are preserved.
\end{lemma}

\begin{lemma}[Success and underflow rewinds in the host]
\lean{Turing.FinTM.loopHost_borrow_rewind}
\label{Turing.FinTM.loopHost_borrow_rewind}
\leanok
\uses{Turing.Action.apply, Turing.Cfg, Turing.Cfg.workTapeSymbols, Turing.FinTM, Turing.FinTM.LoopHostState, Turing.FinTM.bufferTape, Turing.FinTM.bufferTape_left, Turing.FinTM.bufferTape_nat, Turing.FinTM.loopControlAction, Turing.FinTM.loopControl_apply, Turing.FinTM.loopFrame, Turing.FinTM.loopFrame_counter, Turing.FinTM.loopHost, Turing.FinTM.loopWrite, Turing.MultiTapeTM.runFrom, Turing.MultiTapeTM.runFrom_succ_eq_step, Turing.MultiTapeTM.runFrom_zero, Turing.MultiTapeTM.step}
\difficulty{5}
Let $w$ be a counter word, $b$ a success bit, and $j \le |w|$. Running the loop host for
exactly $j + 1$ steps from the rewind phase selected by $b$ on a frame whose counter
track is the buffer tape of $w$ with head at cell $j - 1$ returns the counter head to
zero and dispatches: success releases the next anchor call, underflow enters the
exhaustion phase without yet emitting. All inactive residue is retained.
\end{lemma}

\begin{lemma}[Complete counter operation in the host]
\lean{Turing.FinTM.loopHost_borrow}
\label{Turing.FinTM.loopHost_borrow}
\leanok
\uses{Turing.Cfg, Turing.FinTM, Turing.FinTM.LoopHostState, Turing.FinTM.bufferTape, Turing.FinTM.loopBorrowPos, Turing.FinTM.loopBorrowPos_le, Turing.FinTM.loopDebit, Turing.FinTM.loopDebit_length, Turing.FinTM.loopFrame, Turing.FinTM.loopHost, Turing.FinTM.loopHost_borrow_rewind, Turing.FinTM.loopHost_borrow_run, Turing.MultiTapeTM.runFrom, Turing.MultiTapeTM.runFrom_add}
\difficulty{4}
Let $w$ be a counter word and $j$ the length of its maximal $\mathrm{false}$ prefix.
Then $2j + 2 \le 2|w| + 2$, and running the loop host for exactly $2j + 2$ steps from
its borrow phase on a frame whose counter track is the buffer tape of $w$ with head at
zero ends with the counter track holding the decremented word, head at zero, in the
released next-anchor call state on success and in the exhaustion phase on underflow.
\end{lemma}

\begin{lemma}[The flag track reads its origin cell]
\lean{Turing.FinTM.loopFrame_flag}
\label{Turing.FinTM.loopFrame_flag}
\leanok
\uses{Turing.Cfg, Turing.Cfg.workTapeSymbols, Turing.FinTM, Turing.FinTM.LoopHostState, Turing.FinTM.loopFrame}
\difficulty{1}
For every controller frame of the loop host with flag tape $f$, the symbol scanned on
the flag track of the frame is exactly $f(0)$, independently of all inactive residue.
\end{lemma}

\begin{lemma}[Clearing the only flag cell blanks the flag tape]
\lean{Turing.FinTM.loopFlag_clear}
\label{Turing.FinTM.loopFlag_clear}
\leanok
\uses{Turing.FinTM.bufferTape, Turing.FinTM.loopWrite}
\difficulty{2}
For every optional flag bit, blanking the origin cell of the tape that stores that bit
at the origin and is blank elsewhere yields the completely blank buffer tape.
\end{lemma}

\begin{lemma}[A rejecting round debits or exhausts]
\lean{Turing.FinTM.loopHost_reject}
\label{Turing.FinTM.loopHost_reject}
\leanok
\uses{Turing.Action.apply, Turing.Cfg, Turing.Cfg.workTapeSymbols, Turing.FinTM, Turing.FinTM.bufferTape, Turing.FinTM.leftCfg, Turing.FinTM.loopBodyCfg, Turing.FinTM.loopBodyPadded, Turing.FinTM.loopBorrowPos, Turing.FinTM.loopBorrowPos_le, Turing.FinTM.loopCall, Turing.FinTM.loopCall_frame, Turing.FinTM.loopCall_reframe, Turing.FinTM.loopControlAction, Turing.FinTM.loopControl_apply, Turing.FinTM.loopDebit, Turing.FinTM.loopFlag_clear, Turing.FinTM.loopFrame, Turing.FinTM.loopFrame_flag, Turing.FinTM.loopHost, Turing.FinTM.loopHost_borrow, Turing.FinTM.loopWrite, Turing.MultiTapeTM.runFrom, Turing.MultiTapeTM.runFrom_succ_eq_step, Turing.MultiTapeTM.runFrom_succ_eq_step', Turing.MultiTapeTM.step, Turing.captureCfg}
\difficulty{5}
Let $c$ be a halted body configuration with empty output, $w$ a counter word, and $f$ a
fuel configuration. Then there is a time $t \le 2|w| + 4$ such that, from the stopped
call on $c$ with stop flag $\mathrm{false}$: if the fixed-width decrement of $w$
succeeds, the loop host at time $t$ is at the released next-anchor call with the
decremented counter word; otherwise the host at time $t$ has halted with output empty in
find mode and the single bit $\mathrm{false}$ in decision mode.
\end{lemma}

\begin{lemma}[Phase twelve rewinds the accepted payload]
\lean{Turing.FinTM.loopHost_payload_rewind}
\label{Turing.FinTM.loopHost_payload_rewind}
\leanok
\uses{Turing.Action.apply, Turing.Cfg, Turing.Cfg.workTapeSymbols, Turing.FinTM, Turing.FinTM.LoopHostState, Turing.FinTM.bufferTape, Turing.FinTM.bufferTape_left, Turing.FinTM.bufferTape_nat, Turing.FinTM.emCallIdleTM, Turing.FinTM.loopControlAction, Turing.FinTM.loopControl_apply, Turing.FinTM.loopFrame, Turing.FinTM.loopFrame_payload, Turing.FinTM.loopHost, Turing.FinTM.loopWrite, Turing.MultiTapeTM.runFrom, Turing.MultiTapeTM.runFrom_succ_eq_step, Turing.MultiTapeTM.runFrom_zero, Turing.MultiTapeTM.step}
\difficulty{5}
Let $w$ be a captured payload word and $j \le |w|$. Running the loop host for exactly $j
+ 1$ steps from its payload-rewind phase on a frame whose payload track is the buffer
tape of $w$ with payload head at cell $j - 1$ reaches the replay phase with the payload
head at cell zero and every other frame component unchanged.
\end{lemma}

\begin{lemma}[Replay of a framed payload]
\lean{Turing.FinTM.loopHost_frame_replay}
\label{Turing.FinTM.loopHost_frame_replay}
\leanok
\uses{Turing.Cfg, Turing.FinTM, Turing.FinTM.LoopHostState, Turing.FinTM.bufferTape, Turing.FinTM.loopFrame, Turing.FinTM.loopHost, Turing.FinTM.loopHost_replay, Turing.FinTM.loopReplayCfg, Turing.FinTM.rightCfg, Turing.MultiTapeTM.runFrom}
\difficulty{4}
Let $w$ be a captured payload word and consider the loop host in its replay phase on a
frame with empty output whose payload track is the buffer tape of $w$ with head at zero.
After exactly $|w| + 1$ steps the host has halted and its output is exactly $w$.
\end{lemma}

\begin{lemma}[An accepting round emits its verdict or payload]
\lean{Turing.FinTM.loopHost_accept}
\label{Turing.FinTM.loopHost_accept}
\leanok
\uses{Turing.Action.apply, Turing.Cfg, Turing.Cfg.workTapeSymbols, Turing.FinTM, Turing.FinTM.bufferTape, Turing.FinTM.leftCfg, Turing.FinTM.loopBodyCfg, Turing.FinTM.loopBodyPadded, Turing.FinTM.loopCall, Turing.FinTM.loopCall_frame, Turing.FinTM.loopControlAction, Turing.FinTM.loopControl_apply, Turing.FinTM.loopFrame, Turing.FinTM.loopFrame_flag, Turing.FinTM.loopHost, Turing.FinTM.loopHost_bound, Turing.FinTM.loopHost_frame_replay, Turing.FinTM.loopHost_payload_rewind, Turing.FinTM.loopWrite, Turing.MultiTapeTM.runFrom, Turing.MultiTapeTM.runFrom_add, Turing.MultiTapeTM.runFrom_succ_eq_step, Turing.MultiTapeTM.step, Turing.captureCfg}
\difficulty{5}
Let $c$ be a halted body configuration, $w$ a counter word, and $f$ a fuel
configuration. Then there is a time $t \le 2\,|o| + 3$, where $o$ is the output of $c$,
such that from the stopped call on $c$ with stop flag $\mathrm{true}$ the loop host at
time $t$ has halted with output $o$ in find mode and output the single bit
$\mathrm{true}$ in decision mode.
\end{lemma}

\begin{lemma}[One complete round of the loop host]
\lean{Turing.FinTM.loopHost_round}
\label{Turing.FinTM.loopHost_round}
\leanok
\uses{Turing.Cfg, Turing.Cfg.ofWords, Turing.FinTM, Turing.FinTM.loopCall, Turing.FinTM.loopDebit, Turing.FinTM.loopHost, Turing.FinTM.loopHost_accept, Turing.FinTM.loopHost_anchor_return, Turing.FinTM.loopHost_bound, Turing.FinTM.loopHost_halt_return, Turing.FinTM.loopHost_reject, Turing.FinTM.loop_first_halt, Turing.FinTM.loop_output_length_le, Turing.MultiTapeTM.runFrom, Turing.MultiTapeTM.runFrom_add, Turing.stateWord}
\difficulty{5}
Let $s$ be a round word, $w$ a counter word, $f$ a fuel configuration, and $t > 0$ a
time such that the body, started at the clean anchor seam carrying $s$, does not revisit
the anchor strictly between times zero and $t$ and at time $t$ either halts with a
declared payload (acceptance) or sits at the clean seam carrying the stepped word. Then
there is a time $v \le 3t + 2|w| + 5$ such that from the released call on the seam of
$s$: acceptance halts the host with the payload (find mode) or the bit $\mathrm{true}$
(decision mode); a rejection with successful decrement reaches the released call on the
next seam with the decremented counter; and a rejection with underflow halts with empty
output (find mode) or the bit $\mathrm{false}$ (decision mode).
\end{lemma}

\begin{definition}[Uniform per-segment coefficient]
\lean{Turing.FinTM.loopHost_bound}
\label{Turing.FinTM.loopHost_bound}
\leanok
The uniform constant $10$; it is the maximum of the audited startup, body, counter, and
dispatch coefficients and serves as the per-segment budget multiplier of the loop host.
\end{definition}

\begin{lemma}[Configuration contracts of the loop controller]
\lean{Turing.FinTM.loopHost_contracts}
\label{Turing.FinTM.loopHost_contracts}
\leanok
\uses{Turing.Cfg, Turing.Cfg.ofWords, Turing.FinTM, Turing.FinTM.ComputesFunInTime, Turing.FinTM.emCallIdleTM, Turing.FinTM.loopCall, Turing.FinTM.loopDebit, Turing.FinTM.loopDebit_iterate_length, Turing.FinTM.loopDebit_iterate_value, Turing.FinTM.loopDebit_success, Turing.FinTM.loopHost, Turing.FinTM.loopHost_bound, Turing.FinTM.loopHost_prepare, Turing.FinTM.loopHost_round, Turing.FinTM.loopHost_start, Turing.FinTM.loop_fuel_width, Turing.FinTM.loop_orbit_inv, Turing.MultiTapeTM.initCfg, Turing.MultiTapeTM.runFrom, Turing.MultiTapeTM.runFrom_add, Turing.stateWord}
\difficulty{6}
Let a loop body with anchor, a fuel machine computing the binary digits of $R(|x|)$ in
time $T$, an invariant preserved from the initial word by the step function, a bounded
anchor-free startup, and a per-round contract (positive duration at most $T(|x|)$, no
early anchor revisit, and accept-with-output or advance-to-stepped-seam) be given,
together with an output mode. Then there is a constant $c$ such that for every input $x$
the loop host has a round-configuration family: a startup segment of at most $c (T(|x|)
+ 1)$ steps reaching round zero, empty output at every round configuration up to
$R(|x|)$, a halted terminal at index $R(|x|) + 1$ with output empty in find mode and the
bit $\mathrm{false}$ in decision mode, and for each round $i \le R(|x|)$ a segment of at
most $c (T(|x|) + 1)$ steps that halts with the declared output when round $i$ accepts
and reaches round $i + 1$ otherwise.
\end{lemma}

\begin{theorem}[The configuration-level loop combinator]
\lean{Turing.FinTM.exists_loopCfgTM}
\label{Turing.FinTM.exists_loopCfgTM}
\leanok
\uses{Turing.Cfg, Turing.Cfg.ofWords, Turing.FinTM, Turing.FinTM.ComputesFunInTime, Turing.FinTM.emCallIdleTM, Turing.FinTM.loopHost, Turing.FinTM.loopHost_contracts, Turing.MultiTapeTM.initCfg, Turing.MultiTapeTM.runFrom, Turing.stateWord}
\difficulty{2}
Let a loop body with anchor, a fuel machine computing the binary digits of $R(|x|)$ in
time $T$, an invariant preserved from the initial word by the step function, a bounded
anchor-free startup, and a per-round accept-or-advance contract with verdict output the
single bit $\mathrm{true}$ be given. Then there are a finite machine $E$ and a constant
$c$ such that on every input $x$, $E$ has a family of round configurations: a startup of
at most $c (T(|x|) + 1)$ steps reaching round zero, empty output at rounds $0, \dots,
R(|x|)$, a halted terminal at index $R(|x|) + 1$ with output the single bit
$\mathrm{false}$, and per-round segments of at most $c (T(|x|) + 1)$ steps that halt
with $\mathrm{true}$ on acceptance and advance to the next round configuration
otherwise.
\end{theorem}

\begin{lemma}[Summation with an already-halted terminal]
\lean{Turing.FinTM.loop_halted_run}
\label{Turing.FinTM.loop_halted_run}
\leanok
\uses{Turing.Cfg, Turing.FinTM.emCallIdleTM, Turing.MultiTapeTM, Turing.MultiTapeTM.runFrom, Turing.MultiTapeTM.runFrom_add}
\difficulty{5}
Let $M$ be a multi-tape Turing machine over the Boolean alphabet, let $c_0, \dots, c_N$
be configurations of $M$ on a common input with $c_N$ halted with output the single bit
$\mathrm{false}$, and let $a$ be a Boolean predicate on round indices. Assume for each
$j < N$ that within $B$ steps the machine either halts from $c_j$ with output the single
bit $\mathrm{true}$ (when $a(j)$ holds) or reaches $c_{j+1}$ (when $a(j)$ fails). Then
from $c_0$ the machine halts within $N \cdot B$ steps with output the single bit
recording whether $a(j)$ holds for some $j < N$.
\end{lemma}

\begin{theorem}[The decision loop combinator]
\lean{Turing.FinTM.exists_loopTM}
\label{Turing.FinTM.exists_loopTM}
\leanok
\uses{Turing.Cfg.ofWords, Turing.FinTM, Turing.FinTM.ComputesFunInTime, Turing.FinTM.ComputesInTime, Turing.FinTM.ComputesInTime.mono, Turing.FinTM.emCallIdleTM, Turing.FinTM.exists_loopCfgTM, Turing.FinTM.loop_halted_run, Turing.MultiTapeTM.initCfg, Turing.MultiTapeTM.runFrom, Turing.MultiTapeTM.runFrom_add, Turing.stateWord}
\difficulty{4}
Let a loop body with anchor be given, a fuel machine computing the binary digits of
$R(|x|)$ in time $T$, an invariant that holds at the initial word and is preserved by
the step function, a startup reaching the initial seam within $T(|x|)$ without visiting
the anchor earlier, and a per-round contract: on every admissible word, in positive time
at most $T(|x|)$ with no early anchor revisit, the body either halts with output the
single bit $\mathrm{true}$ or advances to the seam of the stepped word. Then there are a
finite machine $E$ and a constant $c$ such that $E$ computes, within $c (T(n) + 1)(R(n)
+ 2)$ steps on inputs of length $n$, the one-bit answer to whether some orbit point
among the first $R(|x|) + 1$ iterates of the initial word is accepted.
\end{theorem}

\begin{lemma}[Summation selecting the first accepting payload]
\lean{Turing.FinTM.loop_find_run}
\label{Turing.FinTM.loop_find_run}
\leanok
\uses{Turing.Cfg, Turing.FinTM.emCallIdleTM, Turing.MultiTapeTM, Turing.MultiTapeTM.runFrom, Turing.MultiTapeTM.runFrom_add}
\difficulty{5}
Let $M$ be a multi-tape Turing machine over the Boolean alphabet, let $c_0, \dots, c_N$
be configurations of $M$ on a common input with $c_N$ halted with empty output, let $a$
be a Boolean predicate on round indices, and let $w$ assign a payload word to each
index. Assume for each $j < N$ that within $B$ steps the machine either halts from $c_j$
with output $w(j)$ (when $a(j)$ holds) or reaches $c_{j+1}$ (when $a(j)$ fails). Then
from $c_0$ the machine halts within $N \cdot B$ steps, with output $w(i)$ for the least
accepting index $i < N$ when one exists and empty output otherwise.
\end{lemma}

\begin{theorem}[The result-bearing loop combinator]
\lean{Turing.FinTM.exists_loopFindTM}
\label{Turing.FinTM.exists_loopFindTM}
\leanok
\uses{Turing.Cfg.ofWords, Turing.FinTM, Turing.FinTM.ComputesFunInTime, Turing.FinTM.ComputesInTime, Turing.FinTM.ComputesInTime.mono, Turing.FinTM.emCallIdleTM, Turing.FinTM.loopHost, Turing.FinTM.loopHost_contracts, Turing.FinTM.loop_find_run, Turing.MultiTapeTM.initCfg, Turing.MultiTapeTM.runFrom, Turing.MultiTapeTM.runFrom_add, Turing.stateWord}
\difficulty{4}
Let a loop body with anchor be given, a fuel machine computing the binary digits of
$R(|x|)$ in time $T$, an invariant that holds at the initial word and is preserved by
the step function, a bounded anchor-free startup, and a per-round contract in which an
accepting round halts with a declared payload word and a rejecting round advances to the
seam of the stepped word. Then there are a finite machine $E$ and a constant $c$ such
that $E$ computes, within $c (T(n) + 1)(R(n) + 2)$ steps on inputs of length $n$, the
payload of the first accepted orbit point among the first $R(|x|) + 1$ iterates of the
initial word, and the empty word when no orbit point is accepted.
\end{theorem}

\begin{definition}[Zero-tape placeholder machine]
\lean{Turing.FinTM.emCallIdleTM}
\label{Turing.FinTM.emCallIdleTM}
\leanok
\uses{Turing.FinTM, Turing.FinTM.controlAction}
A bundled finite machine over the Boolean alphabet with no work tapes and a single
control state whose only transition is the stationary silent self-loop. It supplies the
empty inactive bank of the virtual-input wrapper; its own transition is never entered in
that use.
\end{definition}

\begin{definition}[Capture-wrapped candidate evaluator]
\lean{Turing.FinTM.emCallEvalTM}
\label{Turing.FinTM.emCallEvalTM}
\leanok
\uses{Turing.FinTM, Turing.FinTM.bufferedCompTM, Turing.FinTM.controlAction, Turing.FinTM.emCallIdleTM, Turing.captureAction}
For a bundled finite machine $M$, the evaluator obtained by wrapping the buffered
virtual-input simulator of $M$ (reading its argument from a designated work tape rather
than the native input) with the capture transformation: every emission of the simulated
run lands on a final capture tape, the completed state is a live self-looping return
state, and no evaluated bit reaches the physical output.
\end{definition}

\begin{definition}[Exact evaluator configuration]
\lean{Turing.FinTM.emCallEvalCfg}
\label{Turing.FinTM.emCallEvalCfg}
\leanok
\uses{Turing.Cfg, Turing.FinTM, Turing.FinTM.bufferedSecondCfg, Turing.FinTM.emCallEvalTM, Turing.FinTM.emCallIdleTM, Turing.captureCfg}
For a configuration $c$ of a machine $M$ on a simulated argument, a virtual boundary
tag, and a physical input-head position, the configuration of the capture-wrapped
evaluator in which the preserved argument occupies tape zero, the simulated work bank of
$c$ follows it, the final tape captures every emission of $c$, and the physical input
head is fixed at the given position.
\end{definition}

\begin{lemma}[Capture commutes with virtual simulation]
\lean{Turing.FinTM.emCall_eval_run}
\label{Turing.FinTM.emCall_eval_run}
\leanok
\uses{Turing.Cfg, Turing.Cfg.Halted, Turing.FinTM, Turing.FinTM.VirtualTag, Turing.FinTM.bufferedCompTM, Turing.FinTM.bufferedSecondCfg, Turing.FinTM.bufferedSecondCfg_run, Turing.FinTM.emCallEvalCfg, Turing.FinTM.emCallEvalTM, Turing.FinTM.emCallIdleTM, Turing.MultiTapeTM.runFrom, Turing.capture_run}
\difficulty{4}
Let $M$ be a bundled finite machine, $c$ a configuration of $M$ on a simulated argument
carrying a valid virtual boundary tag, and $t$ a time before which the run of $M$ from
$c$ stays live. Then there is a valid tag for the endpoint such that running the
capture-wrapped evaluator for $t$ steps from the evaluator configuration of $c$ equals
the evaluator configuration of the run of $M$ from $c$ at time $t$.
\end{lemma}

\begin{lemma}[The prepared candidate is the clean entry seam]
\lean{Turing.FinTM.emCall_eval_initial}
\label{Turing.FinTM.emCall_eval_initial}
\leanok
\uses{Turing.Cfg.init, Turing.Cfg.ofWords, Turing.FinTM, Turing.FinTM.bufferTape, Turing.FinTM.bufferedSecondCfg, Turing.FinTM.emCallEvalCfg, Turing.FinTM.emCallEvalTM, Turing.FinTM.emCallIdleTM, Turing.FinTM.tapeBlocks, Turing.MultiTapeTM.initCfg, Turing.captureCfg, Turing.stateWord}
\difficulty{4}
For every physical input $w$ and argument word $s$, the evaluator configuration built
from the initial configuration of $M$ on $s$, with initial tag and physical input head
at one, equals the canonical seam configuration of the capture-wrapped evaluator whose
first work tape holds $s$, all other work tapes blank, all heads at their origins, and
output empty.
\end{lemma}

\begin{lemma}[The evaluator reaches the actual first halt]
\lean{Turing.FinTM.emCall_eval_first}
\label{Turing.FinTM.emCall_eval_first}
\leanok
\uses{Turing.Cfg.init, Turing.FinTM, Turing.FinTM.ComputesInTime, Turing.FinTM.ComputesInTime.output_unique, Turing.FinTM.VirtualTag, Turing.FinTM.bufferedSecondCfg, Turing.FinTM.computesInTime_iff, Turing.FinTM.emCallEvalCfg, Turing.FinTM.emCallEvalTM, Turing.FinTM.emCall_eval_run, Turing.MultiTapeTM.initCfg, Turing.MultiTapeTM.runFrom, Turing.captureCfg}
\difficulty{5}
Let $M$ compute output $o$ on argument $s$ within time $T$. Then there is a time $t \le
T$ and a boundary tag such that the run of $M$ from its initial configuration on $s$ is
halted at $t$ with output $o$, the capture-wrapped evaluator started on the prepared
candidate does not visit its return state before $t$, and at time $t$ it is exactly the
evaluator configuration of that halted run.
\end{lemma}

\begin{definition}[Half-open visited-interval marker]
\lean{Turing.FinTM.emCallInterval}
\label{Turing.FinTM.emCallInterval}
\leanok
For an integer $\ell$ and a width $m$, the marker tape that holds $\mathrm{true}$
exactly on the cells $z$ with $\ell \le z < \ell + m$ and is blank elsewhere.
\end{definition}

\begin{definition}[Partially erased interval data]
\lean{Turing.FinTM.emCallCleared}
\label{Turing.FinTM.emCallCleared}
\leanok
For a data tape $d$, an integer $\ell$, and a count $j$, the tape that is blank on the
cells $z$ with $\ell \le z < \ell + j$ and agrees with $d$ everywhere else: the contents
after erasing the first $j$ cells of the interval.
\end{definition}

\begin{lemma}[One erasure extends the cleared interval]
\lean{Turing.FinTM.emCall_cleared_step}
\label{Turing.FinTM.emCall_cleared_step}
\leanok
\uses{Turing.FinTM.emCallClearTM, Turing.FinTM.emCallCleared}
\difficulty{3}
For every data tape $d$, integer $\ell$, and count $j$, blanking the cell $\ell + j$ of
the tape obtained from $d$ by erasing the first $j$ interval cells yields the tape
obtained from $d$ by erasing the first $j + 1$ interval cells.
\end{lemma}

\begin{definition}[Three-tape interval cleaner]
\lean{Turing.FinTM.emCallClearTM}
\label{Turing.FinTM.emCallClearTM}
\leanok
\uses{Turing.FinTM, Turing.FinTM.controlAction}
A bundled finite machine over the Boolean alphabet with three work tapes -- data, an
interval marker, and an origin marker -- and four control states, that scans left to the
marked interval's left end, erases the data and the marker in one rightward pass,
rewinds on the surviving origin marker, and finishes in a silent absorbing return state
with all three heads at the origin. All three heads move together, and only observed
tape symbols dispatch control.
\end{definition}

\begin{definition}[Configuration of the interval cleaner]
\lean{Turing.FinTM.emCallClearCfg}
\label{Turing.FinTM.emCallClearCfg}
\leanok
\uses{Turing.Cfg, Turing.FinTM.emCallClearTM}
For an input word, a physical input-head position, a control state, three tape functions
for data, interval marks, and the origin marker, and a common head position $h$, the
configuration of the three-tape interval cleaner whose tapes are as given, whose three
work heads are all at $h$, and whose output is empty.
\end{definition}

\begin{lemma}[The initial left scan finds the interval end]
\lean{Turing.FinTM.emCall_clear_left}
\label{Turing.FinTM.emCall_clear_left}
\leanok
\uses{Turing.Action.apply, Turing.Cfg.workTapeSymbols, Turing.FinTM.bufferTape, Turing.FinTM.emCallClearCfg, Turing.FinTM.emCallClearTM, Turing.FinTM.emCallInterval, Turing.MultiTapeTM.runFrom, Turing.MultiTapeTM.runFrom_succ_eq_step, Turing.MultiTapeTM.step, Turing.moveInputPos_zero}
\difficulty{5}
Let $d$ be a data tape and let an interval of width $m$ start at $\ell$. For every $j <
m$, running the interval cleaner for exactly $j + 2$ steps from its initial state with
all heads at $\ell + j$ reaches its rightward-erasure state with all heads at $\ell$,
independently of blank holes in $d$.
\end{lemma}

\begin{lemma}[Zero erasures change nothing]
\lean{Turing.FinTM.emCall_cleared_zero}
\label{Turing.FinTM.emCall_cleared_zero}
\leanok
\uses{Turing.FinTM.emCallCleared}
\difficulty{1}
For every data tape $d$ and integer $\ell$, erasing the first zero cells of the interval
starting at $\ell$ yields $d$ unchanged.
\end{lemma}

\begin{lemma}[Erasing the whole interval blanks the tape]
\lean{Turing.FinTM.emCall_cleared_all}
\label{Turing.FinTM.emCall_cleared_all}
\leanok
\uses{Turing.FinTM.emCallCleared}
\difficulty{2}
Let $d$ be a data tape that is blank outside the interval of width $m$ starting at
$\ell$. Then erasing all $m$ interval cells of $d$ yields the completely blank tape,
with no assumption on the values inside the interval.
\end{lemma}

\begin{lemma}[The rightward pass erases data and marker together]
\lean{Turing.FinTM.emCall_clear_scan}
\label{Turing.FinTM.emCall_clear_scan}
\leanok
\uses{Turing.Action.apply, Turing.Cfg.workTapeSymbols, Turing.FinTM.bufferTape, Turing.FinTM.emCallClearCfg, Turing.FinTM.emCallClearTM, Turing.FinTM.emCallCleared, Turing.FinTM.emCallInterval, Turing.FinTM.emCall_cleared_step, Turing.FinTM.emCall_cleared_zero, Turing.MultiTapeTM.runFrom, Turing.MultiTapeTM.runFrom_succ_eq_step', Turing.MultiTapeTM.step, Turing.moveInputPos_zero}
\difficulty{5}
Let $d$ be a data tape and let an interval of width $m$ start at $\ell$. For every $j
\le m$, running the interval cleaner for exactly $j$ steps from its rightward-erasure
state with all heads at $\ell$ erases the first $j$ cells of both the data and the
marker, moves all heads to $\ell + j$, and leaves the origin marker intact.
\end{lemma}

\begin{lemma}[Removing the origin marker blanks its tape]
\lean{Turing.FinTM.emCall_origin_erase}
\label{Turing.FinTM.emCall_origin_erase}
\leanok
\uses{Turing.FinTM.bufferTape}
\difficulty{3}
Blanking the origin cell of the tape that holds $\mathrm{true}$ exactly at the origin
yields the completely blank tape.
\end{lemma}

\begin{lemma}[The origin marker guides the final rewind]
\lean{Turing.FinTM.emCall_clear_origin}
\label{Turing.FinTM.emCall_clear_origin}
\leanok
\uses{Turing.Action, Turing.Action.apply, Turing.Cfg.workTapeSymbols, Turing.FinTM.bufferTape, Turing.FinTM.bufferTape_nat, Turing.FinTM.emCallClearCfg, Turing.FinTM.emCallClearTM, Turing.FinTM.emCall_origin_erase, Turing.MultiTapeTM.runFrom, Turing.MultiTapeTM.runFrom_succ_eq_step, Turing.MultiTapeTM.step, Turing.moveInputPos_zero}
\difficulty{4}
For every natural number $n$, running the interval cleaner for exactly $n + 1$ steps
from its rewind state with blank data and marker tapes and all heads at $n$ reaches the
silent return state with all three tapes blank and all heads at the origin, the origin
marker being erased by the final transition.
\end{lemma}

\begin{lemma}[Complete native cleanup of a marked interval]
\lean{Turing.FinTM.emCall_clear_run}
\label{Turing.FinTM.emCall_clear_run}
\leanok
\uses{Turing.Action.apply, Turing.Cfg.workTapeSymbols, Turing.FinTM.bufferTape, Turing.FinTM.emCallClearCfg, Turing.FinTM.emCallClearTM, Turing.FinTM.emCallInterval, Turing.FinTM.emCall_clear_left, Turing.FinTM.emCall_clear_origin, Turing.FinTM.emCall_clear_scan, Turing.FinTM.emCall_cleared_all, Turing.MultiTapeTM.runFrom, Turing.MultiTapeTM.runFrom_add, Turing.MultiTapeTM.runFrom_succ_eq_step', Turing.MultiTapeTM.step, Turing.moveInputPos_zero}
\difficulty{5}
Let $d$ be a data tape blank outside an interval of width $m$ starting at $\ell$, with
$\ell \le 0 < \ell + m$, and let the heads start at $\ell + j$ for some $j < m$. Then
there is a positive time $t \le 3m + 4$ after which the interval cleaner, started in its
initial state on $d$, the interval marker, and the origin marker, is in its silent
return state with all three tapes blank and all heads at the origin.
\end{lemma}

\begin{definition}[Closed visited-span marker]
\lean{Turing.FinTM.emCallSpan}
\label{Turing.FinTM.emCallSpan}
\leanok
For integers $\ell \le h$, the marker tape that holds $\mathrm{true}$ exactly on the
cells $z$ with $\ell \le z \le h$ and is blank elsewhere.
\end{definition}

\begin{lemma}[Marking a neighbouring cell extends the span]
\lean{Turing.FinTM.emCall_span_extend}
\label{Turing.FinTM.emCall_span_extend}
\leanok
\uses{Turing.FinTM.emCallSpan}
\difficulty{3}
Let $\ell \le h$ be integers and let $p$ be a cell with $\ell - 1 \le p \le h + 1$.
Writing the mark at $p$ on the closed-span marker for $[\ell, h]$ yields the closed-span
marker for $[\min(\ell,p), \max(h,p)]$.
\end{lemma}

\begin{definition}[Three parallel tape banks]
\lean{Turing.FinTM.emCallSlots}
\label{Turing.FinTM.emCallSlots}
\leanok
Given three families of values indexed by $k$ tapes -- data, visited marks, and origin
marks -- this is the single family indexed by $k + (k + k)$ tapes that lays the three
banks out consecutively.
\end{definition}

\begin{definition}[Tracked evaluator with visited-cell markers]
\lean{Turing.FinTM.emCallTrackTM}
\label{Turing.FinTM.emCallTrackTM}
\leanok
\uses{Turing.FinTM, Turing.FinTM.emCallSlots}
For a bundled finite machine $M$, the machine with three parallel tape banks that
simulates $M$ at two native steps per source step: the first microstep performs the
source action on the data bank while moving the marker heads identically, and the second
stamps the newly reached cells in the visited bank before dispatching, possibly halting.
A single initialization step marks each origin in both marker banks. The physical output
is the source output.
\end{definition}

\begin{definition}[Completed tracked configuration]
\lean{Turing.FinTM.emCallTrackCfg}
\label{Turing.FinTM.emCallTrackCfg}
\leanok
\uses{Turing.Cfg, Turing.FinTM, Turing.FinTM.bufferTape, Turing.FinTM.emCallSlots, Turing.FinTM.emCallSpan, Turing.FinTM.emCallTrackTM}
For a configuration $c$ of the source machine and per-tape visited bounds $\ell$ and
$h$, the tracked-evaluator configuration whose data bank carries $c$, whose visited bank
holds the closed-span markers $[\ell_i, h_i]$, whose origin bank marks each origin, and
whose three per-tape heads agree with the heads of $c$.
\end{definition}

\begin{definition}[Intermediate stamp configuration]
\lean{Turing.FinTM.emCallTrackMid}
\label{Turing.FinTM.emCallTrackMid}
\leanok
\uses{Turing.Cfg, Turing.FinTM, Turing.FinTM.emCallTrackCfg, Turing.FinTM.emCallTrackTM}
For a source configuration $c$ and visited bounds, the tracked-evaluator configuration
that differs from the completed tracked configuration only in its control state, which
is the stamping state storing the successor control of $c$: the data, input, heads, and
emitted output already reflect the performed action while the visited markers are still
the old ones.
\end{definition}

\begin{lemma}[The source microstep of the tracked evaluator]
\lean{Turing.FinTM.emCall_track_action}
\label{Turing.FinTM.emCall_track_action}
\leanok
\uses{Turing.Action.apply, Turing.Cfg, Turing.Cfg.inputSymbol, Turing.Cfg.workTapeSymbols, Turing.FinTM, Turing.FinTM.emCallSlots, Turing.FinTM.emCallTrackCfg, Turing.FinTM.emCallTrackMid, Turing.FinTM.emCallTrackTM, Turing.MultiTapeTM.step}
\difficulty{5}
Let $c$ be a live configuration of the source machine with visited bounds $\ell$ and
$h$. One step of the tracked evaluator from the completed tracked configuration of $c$
equals the intermediate stamp configuration of the one-step successor of $c$ with the
same bounds: the data bank performs the source write and emission, and all three per-
tape heads move by the same source action.
\end{lemma}

\begin{lemma}[The stamping microstep extends the span]
\lean{Turing.FinTM.emCall_track_stamp}
\label{Turing.FinTM.emCall_track_stamp}
\leanok
\uses{Turing.Action.apply, Turing.Cfg, Turing.FinTM, Turing.FinTM.emCallSlots, Turing.FinTM.emCallTrackCfg, Turing.FinTM.emCallTrackMid, Turing.FinTM.emCallTrackTM, Turing.FinTM.emCall_span_extend, Turing.MultiTapeTM.step, Turing.moveInputPos_zero}
\difficulty{4}
Let $c$ be a source configuration with visited bounds $\ell_i \le h_i$ such that each
head of $c$ lies within one cell of its span. One step of the tracked evaluator from the
intermediate stamp configuration of $c$ equals the completed tracked configuration of
$c$ with each span extended to include the current head, after which the stored
successor control (possibly a halt) is dispatched.
\end{lemma}

\begin{definition}[Leftmost visited head position]
\lean{Turing.FinTM.emCallLo}
\label{Turing.FinTM.emCallLo}
\leanok
\uses{Turing.FinTM, Turing.MultiTapeTM.initCfg, Turing.MultiTapeTM.runFrom}
For a machine $M$, an input $x$, a time $t$, and a work tape $i$, the minimum of the
origin and all head positions of tape $i$ along the run of $M$ from its initial
configuration on $x$ through time $t$; it is the left end of the visited interval of
tape $i$.
\end{definition}

\begin{definition}[Rightmost visited head position]
\lean{Turing.FinTM.emCallHi}
\label{Turing.FinTM.emCallHi}
\leanok
\uses{Turing.FinTM, Turing.MultiTapeTM.initCfg, Turing.MultiTapeTM.runFrom}
For a machine $M$, an input $x$, a time $t$, and a work tape $i$, the maximum of the
origin and all head positions of tape $i$ along the run of $M$ from its initial
configuration on $x$ through time $t$; it is the right end of the visited interval of
tape $i$.
\end{definition}

\begin{lemma}[Basic bounds on the visited interval]
\lean{Turing.FinTM.emCall_track_extent}
\label{Turing.FinTM.emCall_track_extent}
\leanok
\uses{Turing.Cfg.init, Turing.FinTM, Turing.FinTM.emCallHi, Turing.FinTM.emCallLo, Turing.MultiTapeTM.initCfg, Turing.MultiTapeTM.runFrom, Turing.MultiTapeTM.runFrom_succ_eq_step', Turing.MultiTapeTM.workTapePos_step_le}
\difficulty{4}
For every machine $M$, input $x$, time $t$, and work tape $i$: the leftmost visited
position is at most zero, the rightmost is at least zero, the current head at time $t$
lies between them, and they lie within $[-t, t]$, so the visited interval has width at
most $2t + 1$.
\end{lemma}

\begin{lemma}[No writes land outside the visited interval]
\lean{Turing.FinTM.emCall_track_support}
\label{Turing.FinTM.emCall_track_support}
\leanok
\uses{Turing.Action.apply, Turing.Cfg.inputSymbol, Turing.Cfg.workTapeSymbols, Turing.FinTM, Turing.FinTM.emCallHi, Turing.FinTM.emCallLo, Turing.FinTM.emCall_track_extent, Turing.MultiTapeTM.initCfg, Turing.MultiTapeTM.runFrom, Turing.MultiTapeTM.runFrom_succ_eq_step', Turing.MultiTapeTM.step}
\difficulty{4}
For every machine $M$, input $x$, time $t$, work tape $i$, and cell $z$ outside the
visited interval of tape $i$ at time $t$, the cell $z$ of tape $i$ after running $M$ for
$t$ steps from its initial configuration on $x$ is blank.
\end{lemma}

\begin{lemma}[The one-point span is the origin marker]
\lean{Turing.FinTM.emCall_span_zero}
\label{Turing.FinTM.emCall_span_zero}
\leanok
\uses{Turing.FinTM.bufferTape, Turing.FinTM.emCallSpan}
\difficulty{3}
The closed-span marker of the one-point interval $[0,0]$ equals the buffer tape holding
the single bit $\mathrm{true}$.
\end{lemma}

\begin{lemma}[Initialization installs the origin markers]
\lean{Turing.FinTM.emCall_track_initial}
\label{Turing.FinTM.emCall_track_initial}
\leanok
\uses{Turing.Action.apply, Turing.Cfg.init, Turing.FinTM, Turing.FinTM.bufferTape_append, Turing.FinTM.bufferTape_nil, Turing.FinTM.emCallSlots, Turing.FinTM.emCallTrackCfg, Turing.FinTM.emCallTrackTM, Turing.FinTM.emCall_span_zero, Turing.MultiTapeTM.initCfg, Turing.MultiTapeTM.runFrom, Turing.MultiTapeTM.step, Turing.moveInputPos_zero}
\difficulty{4}
For every machine $M$ and input $x$, one step of the tracked evaluator from its initial
configuration on $x$ equals the completed tracked configuration of the initial
configuration of $M$ on $x$ with all visited bounds zero: both marker banks hold origin
marks, the data bank is blank, and the output is empty.
\end{lemma}

\begin{lemma}[Exact tracked simulation at two steps per step]
\lean{Turing.FinTM.emCall_track_run}
\label{Turing.FinTM.emCall_track_run}
\leanok
\uses{Turing.FinTM, Turing.FinTM.emCallHi, Turing.FinTM.emCallLo, Turing.FinTM.emCallTrackCfg, Turing.FinTM.emCallTrackTM, Turing.FinTM.emCall_track_action, Turing.FinTM.emCall_track_extent, Turing.FinTM.emCall_track_initial, Turing.FinTM.emCall_track_stamp, Turing.MultiTapeTM.initCfg, Turing.MultiTapeTM.runFrom, Turing.MultiTapeTM.runFrom_add, Turing.MultiTapeTM.runFrom_of_halt, Turing.MultiTapeTM.runFrom_succ_eq_step', Turing.MultiTapeTM.step, Turing.MultiTapeTM.step_of_halt, Turing.MultiTapeTM.workTapePos_step_le}
\difficulty{5}
For every machine $M$, input $x$, and time $t$, running the tracked evaluator for $1 +
2t$ steps from its initial configuration on $x$ equals the completed tracked
configuration of the run of $M$ at time $t$ with the true leftmost and rightmost visited
bounds, including after the source has halted.
\end{lemma}

\begin{lemma}[Tracking costs a factor of two plus one]
\lean{Turing.FinTM.emCall_track_computes}
\label{Turing.FinTM.emCall_track_computes}
\leanok
\uses{Turing.FinTM, Turing.FinTM.ComputesInTime, Turing.FinTM.computesInTime_iff, Turing.FinTM.emCallTrackCfg, Turing.FinTM.emCallTrackTM, Turing.FinTM.emCall_track_run}
\difficulty{2}
Let $M$ compute output $o$ on input $x$ within time $T$. Then the tracked evaluator of
$M$ computes the same output $o$ on $x$ within time $1 + 2T$.
\end{lemma}

\begin{lemma}[Closed spans as half-open intervals]
\lean{Turing.FinTM.emCall_span_interval}
\label{Turing.FinTM.emCall_span_interval}
\leanok
\uses{Turing.FinTM.emCallInterval, Turing.FinTM.emCallSpan}
\difficulty{3}
For all integers $\ell \le h$, the closed-span marker of $[\ell, h]$ equals the half-
open interval marker starting at $\ell$ of width $h - \ell + 1$.
\end{lemma}

\begin{lemma}[Each tracked bank satisfies the cleaner contract]
\lean{Turing.FinTM.emCall_track_clearable}
\label{Turing.FinTM.emCall_track_clearable}
\leanok
\uses{Turing.FinTM, Turing.FinTM.bufferTape, Turing.FinTM.emCallClearCfg, Turing.FinTM.emCallClearTM, Turing.FinTM.emCallFinishTM, Turing.FinTM.emCallHi, Turing.FinTM.emCallLo, Turing.FinTM.emCallSpan, Turing.FinTM.emCall_clear_run, Turing.FinTM.emCall_span_interval, Turing.FinTM.emCall_track_extent, Turing.FinTM.emCall_track_support, Turing.MultiTapeTM.initCfg, Turing.MultiTapeTM.runFrom}
\difficulty{5}
Let $M$ be a machine, $x$ an input, $T$ a time, and $i$ a work tape. Then there is a
positive time $t \le 6T + 7$ after which the interval cleaner, started on the actual
tape $i$ of the run of $M$ at time $T$, the closed-span marker of the visited interval,
and the origin marker, with heads at the actual head position, reaches its return state
with all three tapes blank and all heads at the origin.
\end{lemma}

\begin{lemma}[Cutting a phase at its first return]
\lean{Turing.FinTM.emCall_first_entry}
\label{Turing.FinTM.emCall_first_entry}
\leanok
\uses{Turing.Cfg, Turing.FinTM.emCallRightTM, Turing.MultiTapeTM, Turing.MultiTapeTM.runFrom, Turing.MultiTapeTM.runFrom_add, Turing.MultiTapeTM.step}
\difficulty{4}
Let $M$ be a multi-tape Turing machine over the Boolean alphabet, let a set of stop
control states be absorbing (a step from any configuration in a stop state changes
nothing), and suppose the run from a configuration $c$ reaches a configuration $d$ in a
stop state at time $T$. Then there is a time $t \le T$ such that no stop state is
visited strictly before $t$ and the run from $c$ at time $t$ equals $d$.
\end{lemma}

\begin{lemma}[Cleaner dispatch at its observed first return]
\lean{Turing.FinTM.emCall_clear_first}
\label{Turing.FinTM.emCall_clear_first}
\leanok
\uses{Turing.Action.apply, Turing.FinTM, Turing.FinTM.bufferTape, Turing.FinTM.controlAction, Turing.FinTM.controlAction_apply, Turing.FinTM.emCallClearCfg, Turing.FinTM.emCallClearTM, Turing.FinTM.emCallHi, Turing.FinTM.emCallLo, Turing.FinTM.emCallRightTM, Turing.FinTM.emCallSpan, Turing.FinTM.emCall_first_entry, Turing.FinTM.emCall_track_clearable, Turing.MultiTapeTM.initCfg, Turing.MultiTapeTM.runFrom, Turing.MultiTapeTM.step, Turing.moveInputPos_zero}
\difficulty{4}
Let $M$ be a machine, $x$ an input, $T$ a time, and $i$ a work tape. Then there is a
positive time $t \le 6T + 7$ such that the interval cleaner, started on the actual tape
$i$ of the run of $M$ at time $T$ with its visited-span and origin markers, does not
visit its return state before $t$ and at time $t$ is exactly the blank-endpoint return
configuration with all heads at the origin.
\end{lemma}

\begin{definition}[Right-boundary normalization wrapper]
\lean{Turing.FinTM.emCallRightTM}
\label{Turing.FinTM.emCallRightTM}
\leanok
\uses{Turing.FinTM, Turing.FinTM.controlAction}
For a bundled finite machine $M$, the machine that follows $M$ exactly while it is live
and, after the halting transition of $M$, moves the input head rightwards to the right
boundary blank and then halts silently; the completed source output and work tapes are
untouched by the scan.
\end{definition}

\begin{definition}[Pre-scan embedding of a source configuration]
\lean{Turing.FinTM.emCallRightCfg}
\label{Turing.FinTM.emCallRightCfg}
\leanok
\uses{Turing.Cfg, Turing.FinTM, Turing.FinTM.emCallRightTM}
For a configuration $c$ of the source machine, the configuration of the right-boundary
wrapper with the same tapes, heads, and output, whose control is the embedded live state
of $c$ when $c$ is live and the administrative scan-entry state when $c$ has halted.
\end{definition}

\begin{lemma}[The wrapper follows one live source step]
\lean{Turing.FinTM.emCall_right_step}
\label{Turing.FinTM.emCall_right_step}
\leanok
\uses{Turing.Cfg, Turing.FinTM, Turing.FinTM.emCallRightCfg, Turing.FinTM.emCallRightTM, Turing.MultiTapeTM.step}
\difficulty{3}
Let $c$ be a live configuration of the source machine. One step of the right-boundary
wrapper from the embedding of $c$ equals the embedding of the one-step successor of $c$,
including any halting emission; only the successor control encoding differs.
\end{lemma}

\begin{lemma}[Exact simulation up to the source halt]
\lean{Turing.FinTM.emCall_right_run}
\label{Turing.FinTM.emCall_right_run}
\leanok
\uses{Turing.Cfg, Turing.FinTM, Turing.FinTM.emCallRightCfg, Turing.FinTM.emCallRightTM, Turing.FinTM.emCall_right_step, Turing.MultiTapeTM.runFrom, Turing.MultiTapeTM.runFrom_succ_eq_step'}
\difficulty{3}
Let $c$ be a configuration of the source machine and $t$ a time before which the source
run from $c$ stays live. Then the right-boundary wrapper run for $t$ steps from the
embedding of $c$ equals the embedding of the source run at time $t$.
\end{lemma}

\begin{definition}[Rightward scan configuration]
\lean{Turing.FinTM.emCallRightScan}
\label{Turing.FinTM.emCallRightScan}
\leanok
\uses{Turing.Cfg, Turing.FinTM, Turing.FinTM.emCallRightTM}
For a completed source configuration $c$, an optional administrative control bit, and a
position $j$ at most the input length, the wrapper configuration whose work tapes,
heads, and output are those of $c$ verbatim and whose input head is at position $j + 1$;
only the input position moves during the scan.
\end{definition}

\begin{lemma}[The scan reaches the right boundary blank]
\lean{Turing.FinTM.emCall_right_scan}
\label{Turing.FinTM.emCall_right_scan}
\leanok
\uses{Turing.Action.apply, Turing.Cfg, Turing.Cfg.inputSymbol, Turing.FinTM, Turing.FinTM.controlAction, Turing.FinTM.controlAction_apply, Turing.FinTM.emCallRightScan, Turing.FinTM.emCallRightTM, Turing.MultiTapeTM.runFrom, Turing.MultiTapeTM.runFrom_succ_eq_step, Turing.MultiTapeTM.runFrom_zero, Turing.MultiTapeTM.step, Turing.inputSymbolInner, Turing.moveInputPos_pos_of_ne_right, Turing.moveInputPos_zero}
\difficulty{5}
Let $c$ be a completed source configuration on an input $x$ and let $j$ and $r$ satisfy
$j + r = |x|$. Running the right-boundary wrapper for exactly $r + 1$ steps from the
live scan configuration at position $j$ halts it silently at the right-boundary scan
configuration at position $|x|$, never confusing a blank work cell with the input end.
\end{lemma}

\begin{lemma}[Entering and completing the right scan]
\lean{Turing.FinTM.emCall_right_finish}
\label{Turing.FinTM.emCall_right_finish}
\leanok
\uses{Turing.Action.apply, Turing.Cfg, Turing.FinTM, Turing.FinTM.controlAction, Turing.FinTM.controlAction_apply, Turing.FinTM.emCallRightCfg, Turing.FinTM.emCallRightScan, Turing.FinTM.emCallRightTM, Turing.FinTM.emCall_right_scan, Turing.MultiTapeTM.runFrom, Turing.MultiTapeTM.runFrom_add, Turing.MultiTapeTM.step, Turing.moveInputPos}
\difficulty{4}
Let $c$ be a halted configuration of the source machine on an input $x$. Then there is a
time $t \le |x| + 2$ such that the right-boundary wrapper, run for $t$ steps from the
embedding of $c$, is the halted scan configuration at position $|x|$, with the work
tapes and output of $c$ unchanged.
\end{lemma}

\begin{lemma}[Endpoint of right-boundary normalization]
\lean{Turing.FinTM.emCall_right_endpoint}
\label{Turing.FinTM.emCall_right_endpoint}
\leanok
\uses{Turing.FinTM, Turing.FinTM.emCallRightCfg, Turing.FinTM.emCallRightScan, Turing.FinTM.emCallRightTM, Turing.FinTM.emCall_right_finish, Turing.FinTM.emCall_right_run, Turing.MultiTapeTM.initCfg, Turing.MultiTapeTM.runFrom, Turing.MultiTapeTM.runFrom_add, Turing.MultiTapeTM.runFrom_of_halt}
\difficulty{5}
Let $M$ be a machine and $x$ an input such that the run of $M$ from its initial
configuration is halted at time $T$. Then the right-boundary wrapper of $M$, run for $T
+ |x| + 2$ steps from its initial configuration on $x$, equals the halted scan
configuration built on the run of $M$ at time $T$, with input head at position $|x| +
1$.
\end{lemma}

\begin{lemma}[Normalization preserves timed computation]
\lean{Turing.FinTM.emCall_right_computes}
\label{Turing.FinTM.emCall_right_computes}
\leanok
\uses{Turing.FinTM, Turing.FinTM.ComputesInTime, Turing.FinTM.computesInTime_iff, Turing.FinTM.emCallRightTM, Turing.FinTM.emCall_right_endpoint, Turing.MultiTapeTM.initCfg, Turing.MultiTapeTM.runFrom}
\difficulty{2}
Let $M$ compute output $o$ on input $x$ within time $T$. Then the right-boundary wrapper
of $M$ computes $o$ on $x$ within time $T + |x| + 2$, and after exactly that many steps
its input head is at position $|x| + 1$.
\end{lemma}

\begin{lemma}[First return of the prepared tracked evaluation]
\lean{Turing.FinTM.emCall_prepared_eval_first}
\label{Turing.FinTM.emCall_prepared_eval_first}
\leanok
\uses{Turing.Cfg.init, Turing.FinTM, Turing.FinTM.ComputesInTime, Turing.FinTM.computesInTime_iff, Turing.FinTM.emCallEvalCfg, Turing.FinTM.emCallEvalTM, Turing.FinTM.emCallHi, Turing.FinTM.emCallLo, Turing.FinTM.emCallRightScan, Turing.FinTM.emCallRightTM, Turing.FinTM.emCallTrackCfg, Turing.FinTM.emCallTrackTM, Turing.FinTM.emCall_eval_first, Turing.FinTM.emCall_right_computes, Turing.FinTM.emCall_right_endpoint, Turing.FinTM.emCall_track_computes, Turing.FinTM.emCall_track_run, Turing.MultiTapeTM.initCfg, Turing.MultiTapeTM.runFrom, Turing.MultiTapeTM.runFrom_add, Turing.MultiTapeTM.runFrom_of_halt, Turing.MultiTapeTM.runFrom_zero}
\difficulty{5}
Let $M$ compute output $o$ on argument $s$ within time $T$, and consider the capture-
wrapped evaluator of the right-boundary wrapper of the tracked evaluator of $M$, started
on the prepared candidate seam for $s$. Then there is a positive time $t \le 1 + 2T +
|s| + 2$ such that the evaluator does not visit its return state before $t$ and at time
$t$ is exactly the evaluator configuration of the halted, right-normalized tracked
endpoint: full trace banks, complete capture buffer, and candidate head at the right
boundary.
\end{lemma}

\begin{definition}[Action relocation onto selected host tapes]
\lean{Turing.FinTM.emCallAction}
\label{Turing.FinTM.emCallAction}
\leanok
\uses{Turing.Action}
For a partial selection assigning to each host tape an optional source tape, an
embedding of source control states into host states, and a source action, the host
action whose selected tapes receive the source operations, whose unselected tapes are
stationary, and whose input move, emission, and successor state are those of the source
action under the state embedding.
\end{definition}

\begin{definition}[Configuration relocation onto selected host tapes]
\lean{Turing.FinTM.emCallCfg}
\label{Turing.FinTM.emCallCfg}
\leanok
\uses{Turing.Cfg}
For a partial tape selection, a control-state embedding, default host tapes and heads
for the unselected slots, and a source configuration $c$, the host configuration whose
selected tapes and heads come from $c$, whose unselected tapes and heads are the
defaults, and whose control, input position, and output are those of $c$ under the state
embedding.
\end{definition}

\begin{lemma}[Relocation commutes with one action]
\lean{Turing.FinTM.emCall_apply}
\label{Turing.FinTM.emCall_apply}
\leanok
\uses{Turing.Action, Turing.Action.apply, Turing.Cfg, Turing.FinTM.emCallAction, Turing.FinTM.emCallCfg, Turing.FinTM.emCallFinishTM}
\difficulty{3}
For every partial tape selection, control embedding, default tapes and heads, source
action $a$, and source configuration $c$: applying the relocated action to the relocated
configuration of $c$ equals the relocated configuration of $a$ applied to $c$, including
the write, head motion, and emission; unselected tapes and heads are fixed.
\end{lemma}

\begin{lemma}[Guarded relocation is exact]
\lean{Turing.FinTM.emCall_relocate_run}
\label{Turing.FinTM.emCall_relocate_run}
\leanok
\uses{Turing.Cfg, Turing.Cfg.workTapeSymbols, Turing.FinTM.emCallAction, Turing.FinTM.emCallCfg, Turing.FinTM.emCallFinishTM, Turing.FinTM.emCall_apply, Turing.MultiTapeTM, Turing.MultiTapeTM.runFrom, Turing.MultiTapeTM.runFrom_succ_eq_step', Turing.MultiTapeTM.step, Turing.MultiTapeTM.step_of_halt}
\difficulty{5}
Let a source machine, a host machine, a tape index map with a partial inverse selection,
and a control embedding be given such that on all good source states the host transition
equals the relocated source transition. If every live state visited by the source run
from a configuration $c$ before time $t$ is good, then the host run for $t$ steps from
the relocated configuration of $c$ equals the relocated configuration of the source run
at time $t$.
\end{lemma}

\begin{definition}[Two-tape result finalizer]
\lean{Turing.FinTM.emCallFinishTM}
\label{Turing.FinTM.emCallFinishTM}
\leanok
\uses{Turing.FinTM, Turing.FinTM.controlAction}
A bundled finite machine over the Boolean alphabet with two work tapes (the argument and
the captured result) and six control states that, entering at the known right blanks of
both words, rewinds the argument -- erasing it in install mode, preserving it in emit
mode -- rewinds the capture, transfers the captured word in order (writing it onto the
argument tape in install mode, emitting it to the physical output in emit mode), then
erases the capture and returns both heads to the origin, finishing in an absorbing
silent state.
\end{definition}

\begin{definition}[Configuration of the result finalizer]
\lean{Turing.FinTM.emCallFinishCfg}
\label{Turing.FinTM.emCallFinishCfg}
\leanok
\uses{Turing.Cfg, Turing.FinTM.emCallFinishTM}
For a mode bit, an input word, a physical input-head position, one of six control
states, two tape functions, two head positions, and an output word, the two-tape
finalizer configuration whose argument and capture tracks are the two given tapes with
the given heads and whose control, input position, and output are as given.
\end{definition}

\begin{lemma}[Erasing the last cell of a buffered word]
\lean{Turing.FinTM.emCall_erase_last}
\label{Turing.FinTM.emCall_erase_last}
\leanok
\uses{Turing.FinTM.bufferTape, Turing.FinTM.bufferTape_append, Turing.FinTM.bufferTape_nat}
\difficulty{3}
For every Boolean word $w$ and bit $b$, blanking cell $|w|$ of the buffer tape holding
$w$ followed by $b$ yields the buffer tape of $w$.
\end{lemma}

\begin{lemma}[The argument pass of the finalizer]
\lean{Turing.FinTM.emCall_finish_arg}
\label{Turing.FinTM.emCall_finish_arg}
\leanok
\uses{Turing.Action.apply, Turing.Cfg.workTapeSymbols, Turing.FinTM.bufferTape, Turing.FinTM.bufferTape_nat, Turing.FinTM.emCallFinishCfg, Turing.FinTM.emCallFinishTM, Turing.FinTM.emCall_erase_last, Turing.MultiTapeTM.runFrom, Turing.MultiTapeTM.runFrom_succ_eq_step, Turing.MultiTapeTM.runFrom_zero, Turing.MultiTapeTM.step, Turing.moveInputPos_zero}
\difficulty{5}
Let the finalizer run in a fixed mode on an argument word $a$ and a capture word $c$,
and let $p$ be a prefix of $a$. Starting from its argument-scan state with the argument
head at cell $|p| - 1$ -- the argument tape holding $a$ in emit mode and only $p$ in
install mode -- after exactly $|p| + 1$ steps it is in the capture-rewind state with the
argument tape holding $a$ in emit mode and empty in install mode, the argument head at
the origin, and the capture head one cell left of $|c|$.
\end{lemma}

\begin{lemma}[The capture rewind of the finalizer]
\lean{Turing.FinTM.emCall_finish_rewind}
\label{Turing.FinTM.emCall_finish_rewind}
\leanok
\uses{Turing.Action.apply, Turing.Cfg.workTapeSymbols, Turing.FinTM.bufferTape, Turing.FinTM.bufferTape_nat, Turing.FinTM.emCallFinishCfg, Turing.FinTM.emCallFinishTM, Turing.MultiTapeTM.runFrom, Turing.MultiTapeTM.runFrom_succ_eq_step, Turing.MultiTapeTM.runFrom_zero, Turing.MultiTapeTM.step, Turing.moveInputPos_zero}
\difficulty{4}
Let the finalizer run in a fixed mode with capture word $c$, and let $p$ be a prefix of
$c$. Starting from its capture-rewind state with the capture head at cell $|p| - 1$,
after exactly $|p| + 1$ steps it is in the transfer state with the capture head at the
origin, the capture word intact, and the argument head still at the origin.
\end{lemma}

\begin{lemma}[The transfer pass of the finalizer]
\lean{Turing.FinTM.emCall_finish_transfer}
\label{Turing.FinTM.emCall_finish_transfer}
\leanok
\uses{Turing.Action.apply, Turing.Cfg.workTapeSymbols, Turing.FinTM.bufferTape, Turing.FinTM.bufferTape_append, Turing.FinTM.bufferTape_nat, Turing.FinTM.emCallFinishCfg, Turing.FinTM.emCallFinishTM, Turing.MultiTapeTM.runFrom, Turing.MultiTapeTM.runFrom_succ_eq_step, Turing.MultiTapeTM.runFrom_zero, Turing.MultiTapeTM.step, Turing.moveInputPos_zero}
\difficulty{5}
Let the finalizer run on an argument word $a$ and capture word $c$, and write $c$ as a
processed prefix $p$ followed by a remainder $r$. Starting from its transfer state with
the capture head at $|p|$ -- having already transferred $p$ to the argument tape
(install mode) or the output (emit mode) -- after exactly $|r| + 1$ steps it is in the
erasure state with the full capture transferred: the argument tape holds $c$ in install
mode and $a$ in emit mode, the output holds $c$ exactly in emit mode, and the capture
word itself is still intact.
\end{lemma}

\begin{lemma}[The erasure pass of the finalizer]
\lean{Turing.FinTM.emCall_finish_erase}
\label{Turing.FinTM.emCall_finish_erase}
\leanok
\uses{Turing.Action.apply, Turing.Cfg.workTapeSymbols, Turing.FinTM.bufferTape, Turing.FinTM.bufferTape_nat, Turing.FinTM.bufferTape_nil, Turing.FinTM.emCallFinishCfg, Turing.FinTM.emCallFinishTM, Turing.FinTM.emCallSource, Turing.FinTM.emCall_erase_last, Turing.MultiTapeTM.runFrom, Turing.MultiTapeTM.runFrom_succ_eq_step, Turing.MultiTapeTM.runFrom_zero, Turing.MultiTapeTM.step, Turing.moveInputPos_zero}
\difficulty{4}
Let the finalizer run in a fixed mode with a remaining captured word $p$. Starting from
its erasure state with the capture head at cell $|p| - 1$ (and, in install mode, the
argument head moving along with it), after exactly $|p| + 1$ steps it is in its final
silent state with the capture tape completely blank and both heads at the origin, the
output untouched.
\end{lemma}

\begin{lemma}[Complete finalization in exact time]
\lean{Turing.FinTM.emCall_finish_run}
\label{Turing.FinTM.emCall_finish_run}
\leanok
\uses{Turing.Action.apply, Turing.FinTM.bufferTape, Turing.FinTM.emCallFinishCfg, Turing.FinTM.emCallFinishTM, Turing.FinTM.emCallSource, Turing.FinTM.emCall_finish_arg, Turing.FinTM.emCall_finish_erase, Turing.FinTM.emCall_finish_rewind, Turing.FinTM.emCall_finish_transfer, Turing.MultiTapeTM.runFrom, Turing.MultiTapeTM.runFrom_add, Turing.MultiTapeTM.runFrom_succ_eq_step, Turing.MultiTapeTM.runFrom_zero, Turing.MultiTapeTM.step, Turing.moveInputPos_zero}
\difficulty{5}
Let the finalizer run in a fixed mode on an argument word $a$ and a capture word $c$,
entering at its start state with the argument head at $|a|$ and the capture head at
$|c|$ and empty output. After exactly $|a| + 3|c| + 5$ steps it is in its final silent
state with both heads at the origin and the capture tape blank; the argument tape holds
$c$ in install mode and $a$ in emit mode, and the output is $c$ in emit mode and empty
in install mode. Empty words are included.
\end{lemma}

\begin{definition}[The prepared source phase of the clean call]
\lean{Turing.FinTM.emCallSource}
\label{Turing.FinTM.emCallSource}
\leanok
\uses{Turing.FinTM, Turing.FinTM.emCallEvalTM, Turing.FinTM.emCallRightTM, Turing.FinTM.emCallTrackTM}
For a bundled finite machine $M$, the capture-wrapped evaluator of the right-boundary
wrapper of the tracked evaluator of $M$: one genuine argument tape, three tracked banks,
and one capture tape, even when $M$ has no work tapes.
\end{definition}

\begin{definition}[Control states of the clean-call controller]
\lean{Turing.FinTM.emCallState}
\label{Turing.FinTM.emCallState}
\leanok
\uses{Turing.FinTM, Turing.FinTM.emCallSource}
The finite control type of the clean-call controller is a disjoint sum of the prepared
source phase's states, one cleaner copy per source tape together with a dispatch index,
and the six finalizer states.
\end{definition}

\begin{definition}[Physical index of a tracked triple slot]
\lean{Turing.FinTM.emCallTripleIndex}
\label{Turing.FinTM.emCallTripleIndex}
\leanok
\uses{Turing.FinTM, Turing.FinTM.bufferedCompTM, Turing.FinTM.emCallEvalTM, Turing.FinTM.emCallIdleTM, Turing.FinTM.emCallRightTM, Turing.FinTM.emCallSource, Turing.FinTM.emCallTrackTM}
For a source tape $i$ of the machine $M$ and a component $r$ among data, visited marks,
and origin marks, this is the physical tape index of component $r$ of bank $i$ in the
clean-call layout.
\end{definition}

\begin{definition}[Partial inverse of the triple layout]
\lean{Turing.FinTM.emCallTripleSelect}
\label{Turing.FinTM.emCallTripleSelect}
\leanok
\uses{Turing.FinTM, Turing.FinTM.emCallSource}
For a source tape $i$ of the machine $M$, the partial selection on physical call-layout
tapes that returns the matching component at the data, visited, and origin slots of bank
$i$ and is undefined elsewhere.
\end{definition}

\begin{lemma}[Triple selection inverts the triple index]
\lean{Turing.FinTM.emCall_triple_inverse}
\label{Turing.FinTM.emCall_triple_inverse}
\leanok
\uses{Turing.FinTM, Turing.FinTM.emCallTripleIndex, Turing.FinTM.emCallTripleSelect}
\difficulty{3}
For every source tape $i$ and component $r$ among data, visited marks, and origin marks,
selecting at the physical index of component $r$ of bank $i$ returns exactly $r$.
\end{lemma}

\begin{definition}[Physical indices of argument and capture]
\lean{Turing.FinTM.emCallPairIndex}
\label{Turing.FinTM.emCallPairIndex}
\leanok
\uses{Turing.FinTM, Turing.FinTM.bufferedCompTM, Turing.FinTM.emCallEvalTM, Turing.FinTM.emCallIdleTM, Turing.FinTM.emCallRightTM, Turing.FinTM.emCallSource, Turing.FinTM.emCallTrackTM}
The map sending the finalizer's two tape roles to physical call-layout indices: the
argument to the first physical tape and the capture to the last.
\end{definition}

\begin{definition}[Partial inverse of the pair layout]
\lean{Turing.FinTM.emCallPairSelect}
\label{Turing.FinTM.emCallPairSelect}
\leanok
\uses{Turing.FinTM, Turing.FinTM.bufferedCompTM, Turing.FinTM.emCallIdleTM, Turing.FinTM.emCallRightTM, Turing.FinTM.emCallSource, Turing.FinTM.emCallTrackTM}
The partial selection on physical call-layout tapes that returns the argument role at
the argument slot and the capture role at the capture slot and is undefined elsewhere.
\end{definition}

\begin{lemma}[Pair selection inverts the pair index]
\lean{Turing.FinTM.emCall_pair_inverse}
\label{Turing.FinTM.emCall_pair_inverse}
\leanok
\uses{Turing.FinTM, Turing.FinTM.bufferedCompTM, Turing.FinTM.emCallIdleTM, Turing.FinTM.emCallPairIndex, Turing.FinTM.emCallPairSelect, Turing.FinTM.emCallRightTM, Turing.FinTM.emCallTrackTM}
\difficulty{2}
For each of the finalizer's two tape roles, selecting at its physical call-layout index
returns exactly that role.
\end{lemma}

\begin{definition}[The native clean-call controller]
\lean{Turing.FinTM.emCallTM}
\label{Turing.FinTM.emCallTM}
\leanok
\uses{Turing.FinTM, Turing.FinTM.controlAction, Turing.FinTM.emCallAction, Turing.FinTM.emCallClearTM, Turing.FinTM.emCallFinishTM, Turing.FinTM.emCallPairIndex, Turing.FinTM.emCallPairSelect, Turing.FinTM.emCallRightTM, Turing.FinTM.emCallSource, Turing.FinTM.emCallState, Turing.FinTM.emCallTrackTM, Turing.FinTM.emCallTripleIndex, Turing.FinTM.emCallTripleSelect}
For a bundled finite machine $M$ and a mode bit, the controller machine that evaluates
$M$ on a tape-resident argument through the prepared tracked source and captures the
result at its observed return, then sequentially runs the interval cleaner on each
tracked bank, each dispatched at its observed return, and finally installs the captured
result on the argument tape (install mode) or forwards it to the physical output while
preserving the argument (emit mode), erasing the capture. No deadline or source running
time occurs in its transition table.
\end{definition}

\begin{definition}[Canonical tape layout of the clean call]
\lean{Turing.FinTM.emCallLayout}
\label{Turing.FinTM.emCallLayout}
\leanok
\uses{Turing.FinTM, Turing.FinTM.emCallSlots, Turing.FinTM.emCallSource, Turing.FinTM.tapeBlocks}
For per-slot values -- an argument value, a capture value, and data, visited, and origin
values for each source tape -- the family over all physical call-layout tapes whose
order is: argument, data bank, visited bank, origin bank, capture.
\end{definition}

\begin{lemma}[Each layout index is one of three kinds]
\lean{Turing.FinTM.emCall_layout_cases}
\label{Turing.FinTM.emCall_layout_cases}
\leanok
\uses{Turing.FinTM, Turing.FinTM.bufferedCompTM, Turing.FinTM.emCallEvalTM, Turing.FinTM.emCallIdleTM, Turing.FinTM.emCallPairIndex, Turing.FinTM.emCallRightTM, Turing.FinTM.emCallSource, Turing.FinTM.emCallTrackTM, Turing.FinTM.emCallTripleIndex}
\difficulty{4}
Every physical tape index of the clean-call layout is the argument slot, the capture
slot, or the slot of one component (data, visited, or origin) of one source bank.
\end{lemma}

\begin{lemma}[Projecting a triple slot from the layout]
\lean{Turing.FinTM.emCall_layout_triple}
\label{Turing.FinTM.emCall_layout_triple}
\leanok
\uses{Turing.FinTM, Turing.FinTM.emCallLayout, Turing.FinTM.emCallSlots, Turing.FinTM.emCallTripleIndex, Turing.FinTM.tapeBlocks}
\difficulty{4}
For every assignment of per-slot values and every source tape $i$, the clean-call layout
evaluated at the physical index of component $r$ of bank $i$ gives the data value, the
visited value, or the origin value of tape $i$ according to $r$.
\end{lemma}

\begin{lemma}[Projecting the administrative slots]
\lean{Turing.FinTM.emCall_layout_pair}
\label{Turing.FinTM.emCall_layout_pair}
\leanok
\uses{Turing.FinTM, Turing.FinTM.bufferedCompTM, Turing.FinTM.emCallIdleTM, Turing.FinTM.emCallLayout, Turing.FinTM.emCallPairIndex, Turing.FinTM.emCallRightTM, Turing.FinTM.emCallSlots, Turing.FinTM.emCallTrackTM, Turing.FinTM.tapeBlocks}
\difficulty{2}
For every assignment of per-slot values, the clean-call layout evaluated at the argument
slot gives the argument value and evaluated at the capture slot gives the capture value.
\end{lemma}

\begin{lemma}[Cleaners do not select administrative tapes]
\lean{Turing.FinTM.emCall_triple_pair}
\label{Turing.FinTM.emCall_triple_pair}
\leanok
\uses{Turing.FinTM, Turing.FinTM.bufferedCompTM, Turing.FinTM.emCallIdleTM, Turing.FinTM.emCallPairIndex, Turing.FinTM.emCallRightTM, Turing.FinTM.emCallTrackTM, Turing.FinTM.emCallTripleSelect}
\difficulty{3}
For every source tape $i$, the triple selection of bank $i$ is undefined at the argument
slot and at the capture slot.
\end{lemma}

\begin{lemma}[A cleaner acts only on its own bank]
\lean{Turing.FinTM.emCall_triple_other}
\label{Turing.FinTM.emCall_triple_other}
\leanok
\uses{Turing.FinTM, Turing.FinTM.emCallFrame, Turing.FinTM.emCallTripleIndex, Turing.FinTM.emCallTripleSelect, Turing.FinTM.emCall_triple_inverse}
\difficulty{4}
For all source tapes $i$ and $j$ and every component $r$, the triple selection of bank
$i$ evaluated at the physical index of component $r$ of bank $j$ returns $r$ when $j =
i$ and selects nothing otherwise.
\end{lemma}

\begin{lemma}[The finalizer does not select tracked storage]
\lean{Turing.FinTM.emCall_pair_triple}
\label{Turing.FinTM.emCall_pair_triple}
\leanok
\uses{Turing.FinTM, Turing.FinTM.bufferedCompTM, Turing.FinTM.emCallIdleTM, Turing.FinTM.emCallPairSelect, Turing.FinTM.emCallRightTM, Turing.FinTM.emCallTrackTM, Turing.FinTM.emCallTripleIndex}
\difficulty{3}
For every source tape $i$ and component $r$, the pair selection of the finalizer is
undefined at the physical index of component $r$ of bank $i$.
\end{lemma}

\begin{definition}[Full host frame during cleanup]
\lean{Turing.FinTM.emCallFrame}
\label{Turing.FinTM.emCallFrame}
\leanok
\uses{Turing.Cfg, Turing.FinTM, Turing.FinTM.bufferTape, Turing.FinTM.emCallLayout, Turing.FinTM.emCallState, Turing.FinTM.emCallTM}
For a mode bit, an input word, an argument word $a$, a capture word $c$, a control
state, per-tape data, visited, and origin banks, and per-tape head positions, the clean-
call configuration whose argument track is the buffer tape of $a$ with head at $|a|$,
whose capture track is the buffer tape of $c$ with head at $|c|$, whose three banks are
as given with their common per-tape heads, with input head at one and empty output.
\end{definition}

\begin{definition}[Frame with a cleaned prefix of banks]
\lean{Turing.FinTM.emCallBankFrame}
\label{Turing.FinTM.emCallBankFrame}
\leanok
\uses{Turing.Cfg, Turing.FinTM, Turing.FinTM.bufferTape, Turing.FinTM.emCallFrame, Turing.FinTM.emCallHi, Turing.FinTM.emCallLo, Turing.FinTM.emCallSpan, Turing.FinTM.emCallState, Turing.FinTM.emCallTM, Turing.MultiTapeTM.initCfg, Turing.MultiTapeTM.runFrom}
For a mode bit, input and argument and capture words, a time $T$, a count $j$, and a
control state, the cleanup frame in which the first $j$ source banks are completely
blank with heads at the origin, while every remaining bank holds the exact halted trace
of the source at time $T$: its data tape, the closed-span marker of its visited
interval, the origin marker, and its true head position.
\end{definition}

\begin{lemma}[Cleaner entry reproduces the bank frame]
\lean{Turing.FinTM.emCall_bank_initial}
\label{Turing.FinTM.emCall_bank_initial}
\leanok
\uses{Turing.FinTM, Turing.FinTM.bufferTape, Turing.FinTM.emCallBankFrame, Turing.FinTM.emCallCfg, Turing.FinTM.emCallClearCfg, Turing.FinTM.emCallFrame, Turing.FinTM.emCallHi, Turing.FinTM.emCallLo, Turing.FinTM.emCallSpan, Turing.FinTM.emCallTripleSelect, Turing.FinTM.emCall_layout_cases, Turing.FinTM.emCall_layout_triple, Turing.FinTM.emCall_triple_inverse, Turing.FinTM.emCall_triple_other, Turing.FinTM.emCall_triple_pair, Turing.MultiTapeTM.initCfg, Turing.MultiTapeTM.runFrom}
\difficulty{4}
For every mode, input, argument and capture words, time $T$, and source tape $i$:
relocating the interval cleaner's entry configuration for bank $i$ (actual data, visited
span, origin marker, actual head) into the cleanup frame whose first $i$ banks are blank
yields exactly that frame, with every other bank and both administrative tapes retained.
\end{lemma}

\begin{lemma}[Cleaner endpoint advances the cleaned prefix]
\lean{Turing.FinTM.emCall_bank_final}
\label{Turing.FinTM.emCall_bank_final}
\leanok
\uses{Turing.FinTM, Turing.FinTM.emCallBankFrame, Turing.FinTM.emCallCfg, Turing.FinTM.emCallClearCfg, Turing.FinTM.emCallFrame, Turing.FinTM.emCallTripleSelect, Turing.FinTM.emCall_layout_cases, Turing.FinTM.emCall_layout_pair, Turing.FinTM.emCall_layout_triple, Turing.FinTM.emCall_triple_inverse, Turing.FinTM.emCall_triple_other, Turing.FinTM.emCall_triple_pair}
\difficulty{4}
For every mode, input, argument and capture words, time $T$, and source tape $i$:
relocating the interval cleaner's blank return endpoint into the cleanup frame whose
first $i$ banks are blank yields the cleanup frame whose first $i + 1$ banks are blank,
in the cleaner-return control state for bank $i$.
\end{lemma}

\begin{lemma}[One bank segment within 6T+8 steps]
\lean{Turing.FinTM.emCall_bank_step}
\label{Turing.FinTM.emCall_bank_step}
\leanok
\uses{Turing.FinTM, Turing.FinTM.bufferTape, Turing.FinTM.controlAction_apply, Turing.FinTM.emCallBankFrame, Turing.FinTM.emCallClearCfg, Turing.FinTM.emCallClearTM, Turing.FinTM.emCallFrame, Turing.FinTM.emCallHi, Turing.FinTM.emCallLo, Turing.FinTM.emCallSpan, Turing.FinTM.emCallTM, Turing.FinTM.emCallTripleIndex, Turing.FinTM.emCallTripleSelect, Turing.FinTM.emCall_bank_final, Turing.FinTM.emCall_bank_initial, Turing.FinTM.emCall_clear_first, Turing.FinTM.emCall_relocate_run, Turing.FinTM.emCall_triple_inverse, Turing.MultiTapeTM.initCfg, Turing.MultiTapeTM.runFrom, Turing.MultiTapeTM.runFrom_succ_eq_step', Turing.MultiTapeTM.step, Turing.moveInputPos_zero}
\difficulty{5}
For every mode, input, argument and capture words, time $T$, and source tape $i$, there
is a time $t \le 6T + 8$ such that the clean-call controller, run for $t$ steps from the
cleanup frame with $i$ cleaned banks in the cleaner-entry state for bank $i$, reaches
the cleanup frame with $i + 1$ cleaned banks in the cleaner-entry state for bank $i +
1$; the dispatch happens at the cleaner's observed return.
\end{lemma}

\begin{lemma}[Sequential cleanup of all banks]
\lean{Turing.FinTM.emCall_banks_run}
\label{Turing.FinTM.emCall_banks_run}
\leanok
\uses{Turing.FinTM, Turing.FinTM.emCallBankFrame, Turing.FinTM.emCallTM, Turing.FinTM.emCall_bank_step, Turing.MultiTapeTM.runFrom, Turing.MultiTapeTM.runFrom_add}
\difficulty{4}
For every mode, input, argument and capture words, time $T$, and count $j \le k$ (the
number of source tapes), there is a time $t \le j (6T + 8)$ such that the clean-call
controller, run for $t$ steps from the cleanup frame with no cleaned banks, reaches the
cleanup frame with $j$ cleaned banks, both administrative words preserved. The claim is
a full configuration equality.
\end{lemma}

\begin{lemma}[The evaluator entry is the argument seam]
\lean{Turing.FinTM.emCall_prepare_initial}
\label{Turing.FinTM.emCall_prepare_initial}
\leanok
\uses{Turing.Cfg.ofWords, Turing.FinTM, Turing.FinTM.emCallCfg, Turing.FinTM.emCallEvalCfg, Turing.FinTM.emCallRightTM, Turing.FinTM.emCallSource, Turing.FinTM.emCallState, Turing.FinTM.emCallTM, Turing.FinTM.emCallTrackTM, Turing.FinTM.emCall_eval_initial, Turing.MultiTapeTM.initCfg, Turing.stateWord}
\difficulty{2}
For every mode, input word, and argument word $a$: the identity relocation of the
prepared evaluator's entry configuration on $a$ equals the clean-call controller's
canonical entry seam, whose first work tape holds $a$ with all other tapes blank, heads
at their origins, and empty output.
\end{lemma}

\begin{lemma}[The completed evaluation is the uncleaned frame]
\lean{Turing.FinTM.emCall_prepare_final}
\label{Turing.FinTM.emCall_prepare_final}
\leanok
\uses{Turing.Cfg.init, Turing.FinTM, Turing.FinTM.bufferedSecondCfg, Turing.FinTM.emCallBankFrame, Turing.FinTM.emCallCfg, Turing.FinTM.emCallEvalCfg, Turing.FinTM.emCallFrame, Turing.FinTM.emCallHi, Turing.FinTM.emCallLayout, Turing.FinTM.emCallLo, Turing.FinTM.emCallRightScan, Turing.FinTM.emCallRightTM, Turing.FinTM.emCallSource, Turing.FinTM.emCallState, Turing.FinTM.emCallTrackCfg, Turing.FinTM.emCallTrackTM, Turing.MultiTapeTM.initCfg, Turing.MultiTapeTM.runFrom, Turing.captureCfg}
\difficulty{3}
Let the source run on argument $a$ have output $c$ at time $T$. The identity relocation
of the evaluator configuration of the halted right-normalized tracked endpoint equals
the cleanup frame with zero cleaned banks: complete source data, visited spans, origin
markers, captured result $c$, and the argument head at its right boundary.
\end{lemma}

\begin{lemma}[From the entry seam to the first bank]
\lean{Turing.FinTM.emCall_prepare_run}
\label{Turing.FinTM.emCall_prepare_run}
\leanok
\uses{Turing.Cfg.ofWords, Turing.FinTM, Turing.FinTM.ComputesInTime, Turing.FinTM.computesInTime_iff, Turing.FinTM.controlAction_apply, Turing.FinTM.emCallBankFrame, Turing.FinTM.emCallEvalCfg, Turing.FinTM.emCallFrame, Turing.FinTM.emCallRightTM, Turing.FinTM.emCallSource, Turing.FinTM.emCallState, Turing.FinTM.emCallTM, Turing.FinTM.emCallTrackTM, Turing.FinTM.emCall_prepare_final, Turing.FinTM.emCall_prepare_initial, Turing.FinTM.emCall_prepared_eval_first, Turing.FinTM.emCall_relocate_run, Turing.MultiTapeTM.initCfg, Turing.MultiTapeTM.runFrom, Turing.MultiTapeTM.runFrom_succ_eq_step', Turing.MultiTapeTM.step, Turing.moveInputPos_zero, Turing.stateWord}
\difficulty{5}
Let $M$ compute output $c$ on argument $a$ within time $T$. For every mode and input
word, there is a time $t \le 2T + |a| + 4$ such that the clean-call controller, run for
$t$ steps from its canonical entry seam on $a$, reaches the cleanup frame with zero
cleaned banks in the first cleaner-entry state, with no scan of the native input.
\end{lemma}

\begin{lemma}[The relocated finalizer entry is the clean frame]
\lean{Turing.FinTM.emCall_finish_initial}
\label{Turing.FinTM.emCall_finish_initial}
\leanok
\uses{Turing.FinTM, Turing.FinTM.bufferTape, Turing.FinTM.emCallBankFrame, Turing.FinTM.emCallCfg, Turing.FinTM.emCallFinishCfg, Turing.FinTM.emCallFrame, Turing.FinTM.emCallPairSelect, Turing.FinTM.emCallState, Turing.FinTM.emCall_layout_cases, Turing.FinTM.emCall_layout_pair, Turing.FinTM.emCall_layout_triple, Turing.FinTM.emCall_pair_inverse, Turing.FinTM.emCall_pair_triple}
\difficulty{4}
For every mode, input, argument word $a$, capture word $c$, and time $T$: relocating the
finalizer's start configuration (argument head at $|a|$, capture head at $|c|$, empty
output) through the pair layout equals the cleanup frame in which all source banks are
blank, in the finalizer-entry control state.
\end{lemma}

\begin{lemma}[The finalizer endpoint is the exit seam]
\lean{Turing.FinTM.emCall_finish_final}
\label{Turing.FinTM.emCall_finish_final}
\leanok
\uses{Turing.Cfg.ofWords, Turing.FinTM, Turing.FinTM.bufferTape, Turing.FinTM.bufferedCompTM, Turing.FinTM.emCallCfg, Turing.FinTM.emCallFinishCfg, Turing.FinTM.emCallIdleTM, Turing.FinTM.emCallPairIndex, Turing.FinTM.emCallPairSelect, Turing.FinTM.emCallRightTM, Turing.FinTM.emCallState, Turing.FinTM.emCallTM, Turing.FinTM.emCallTrackTM, Turing.FinTM.emCallTripleIndex, Turing.FinTM.emCall_layout_cases, Turing.FinTM.emCall_pair_inverse, Turing.FinTM.emCall_pair_triple, Turing.stateWord}
\difficulty{4}
For every mode, input, argument word $a$, and capture word $c$: relocating the
finalizer's final configuration through the pair layout equals the canonical seam in the
exit control state whose first work tape holds $c$ in install mode and $a$ in emit mode,
all other tapes blank, all heads at their origins, with output $c$ in emit mode and
empty otherwise.
\end{lemma}

\begin{lemma}[Finalization inside the full controller]
\lean{Turing.FinTM.emCall_finalize_run}
\label{Turing.FinTM.emCall_finalize_run}
\leanok
\uses{Turing.Cfg.ofWords, Turing.FinTM, Turing.FinTM.bufferTape, Turing.FinTM.controlAction_apply, Turing.FinTM.emCallBankFrame, Turing.FinTM.emCallFinishCfg, Turing.FinTM.emCallFinishTM, Turing.FinTM.emCallFrame, Turing.FinTM.emCallPairIndex, Turing.FinTM.emCallPairSelect, Turing.FinTM.emCallState, Turing.FinTM.emCallTM, Turing.FinTM.emCall_finish_final, Turing.FinTM.emCall_finish_initial, Turing.FinTM.emCall_finish_run, Turing.FinTM.emCall_pair_inverse, Turing.FinTM.emCall_relocate_run, Turing.MultiTapeTM.runFrom, Turing.MultiTapeTM.runFrom_succ_eq_step, Turing.MultiTapeTM.step, Turing.moveInputPos_zero, Turing.stateWord}
\difficulty{4}
For every mode, input, argument word $a$, capture word $c$, and time $T$: running the
clean-call controller for exactly $|a| + 3|c| + 6$ steps from the cleanup frame in which
all banks are clean (in the last cleaner-dispatch state) reaches the canonical exit seam
carrying $c$ in install mode and $a$ in emit mode, with output $c$ exactly in emit mode,
every scratch tape blank, and every head at its origin.
\end{lemma}

\begin{lemma}[The complete clean call within a linear budget]
\lean{Turing.FinTM.emCall_complete}
\label{Turing.FinTM.emCall_complete}
\leanok
\uses{Turing.Cfg.ofWords, Turing.FinTM, Turing.FinTM.ComputesInTime, Turing.FinTM.emCallState, Turing.FinTM.emCallTM, Turing.FinTM.emCall_banks_run, Turing.FinTM.emCall_finalize_run, Turing.FinTM.emCall_prepare_run, Turing.MultiTapeTM.runFrom, Turing.MultiTapeTM.runFrom_add, Turing.stateWord}
\difficulty{4}
Let $M$ be a machine with $k$ work tapes computing output $c$ on argument $a$ within
time $T$. For every mode and input word, there is a time $t \le (24 + 13k)(T + |a| + |c|
+ 1)$ such that the clean-call controller, run for $t$ steps from its canonical entry
seam on $a$, reaches the canonical exit seam carrying $c$ in install mode and $a$ in
emit mode, with output $c$ exactly in emit mode.
\end{lemma}

\begin{lemma}[The exit state is absorbing]
\lean{Turing.FinTM.emCall_exit_fixed}
\label{Turing.FinTM.emCall_exit_fixed}
\leanok
\uses{Turing.Cfg, Turing.FinTM, Turing.FinTM.controlAction, Turing.FinTM.controlAction_apply, Turing.FinTM.emCallAction, Turing.FinTM.emCallFinishTM, Turing.FinTM.emCallPairSelect, Turing.FinTM.emCallTM, Turing.MultiTapeTM.step, Turing.moveInputPos_zero}
\difficulty{3}
For every mode and every configuration $z$ of the clean-call controller whose control is
the designated exit state, one step of the controller from $z$ equals $z$: the exit
state is silent and absorbing on every tape configuration.
\end{lemma}

\begin{lemma}[The clean call returns at its first positive exit]
\lean{Turing.FinTM.emCall_first}
\label{Turing.FinTM.emCall_first}
\leanok
\uses{Turing.Cfg.ofWords, Turing.FinTM, Turing.FinTM.ComputesInTime, Turing.FinTM.emCallState, Turing.FinTM.emCallTM, Turing.FinTM.emCall_complete, Turing.FinTM.emCall_exit_fixed, Turing.FinTM.emCall_first_entry, Turing.MultiTapeTM.runFrom, Turing.stateWord}
\difficulty{4}
Let $M$ be a machine with $k$ work tapes computing output $c$ on argument $a$ within
time $T$. For every mode and input word, there is a positive time $t \le (24 + 13k)(T +
|a| + |c| + 1)$ such that the clean-call controller, started at its canonical entry seam
on $a$, does not visit the exit state before $t$ and at time $t$ is exactly the
canonical exit seam carrying $c$ in install mode and $a$ in emit mode, with output $c$
exactly in emit mode.
\end{lemma}

\begin{lemma}[One step commutes with an output prefix]
\lean{Turing.FinTM.emLoop_step_prefix}
\label{Turing.FinTM.emLoop_step_prefix}
\leanok
\uses{Turing.Action.apply, Turing.Cfg, Turing.Cfg.inputSymbol, Turing.Cfg.workTapeSymbols, Turing.MultiTapeTM, Turing.MultiTapeTM.step}
\difficulty{3}
Let $M$ be a multi-tape Turing machine over the Boolean alphabet, $c$ a configuration,
and $p$ a word. One step of $M$ from $c$ with $p$ prepended to its output equals the
one-step successor of $c$ with $p$ prepended to its output: transitions never inspect
the append-only output.
\end{lemma}

\begin{lemma}[Runs commute with an output prefix]
\lean{Turing.FinTM.emLoop_run_prefix}
\label{Turing.FinTM.emLoop_run_prefix}
\leanok
\uses{Turing.Cfg, Turing.FinTM.emLoop_step_prefix, Turing.MultiTapeTM, Turing.MultiTapeTM.runFrom, Turing.MultiTapeTM.runFrom_succ_eq_step'}
\difficulty{3}
Let $M$ be a multi-tape Turing machine over the Boolean alphabet, $c$ a configuration,
$p$ a word, and $t$ a time. Running $M$ for $t$ steps from $c$ with $p$ prepended to its
output equals the run of $M$ from $c$ at time $t$ with $p$ prepended to its output,
including halted tails.
\end{lemma}

\begin{definition}[The forwarding loop controller]
\lean{Turing.FinTM.emLoopHost}
\label{Turing.FinTM.emLoopHost}
\leanok
\uses{Turing.FinTM, Turing.FinTM.LoopHostState, Turing.FinTM.leftAction, Turing.FinTM.loopBodySource, Turing.FinTM.loopHost, Turing.emitAction}
Given a loop body with anchor, a fuel machine, and a mode bit, the variant of the loop
controller whose body calls route every body emission directly to the physical output
(instead of capturing it), while the fuel capture, counter countdown, and all
administrative transitions are those of the original controller.
\end{definition}

\begin{lemma}[Fuel capture in the forwarding host]
\lean{Turing.FinTM.emLoopHost_fuel_capture}
\label{Turing.FinTM.emLoopHost_fuel_capture}
\leanok
\uses{Turing.Cfg, Turing.Cfg.Halted, Turing.FinTM, Turing.FinTM.emLoopHost, Turing.FinTM.loopFuelSource, Turing.MultiTapeTM.runFrom, Turing.captureCfg, Turing.capture_run}
\difficulty{1}
Let $c$ be a configuration of the relocated fuel source on the forwarding host's tape
layout and $t$ a time before which that source does not halt. Then the forwarding host
run for $t$ steps from the capture embedding of $c$ equals the capture embedding of the
fuel source run at time $t$.
\end{lemma}

\begin{lemma}[Initial configuration of the forwarding host]
\lean{Turing.FinTM.emLoopHost_init}
\label{Turing.FinTM.emLoopHost_init}
\leanok
\uses{Turing.Cfg.ofWords, Turing.FinTM, Turing.FinTM.emLoopHost, Turing.FinTM.loopFuelSource, Turing.FinTM.loopHost, Turing.MultiTapeTM.initCfg, Turing.captureCfg, Turing.initCfg_ofWords}
\difficulty{2}
For every input word $x$, the initial configuration of the forwarding host on $x$ equals
the capture embedding (with empty buffers and return phase zero) of the initial
configuration of the relocated fuel source on $x$.
\end{lemma}

\begin{lemma}[Input rewind in the forwarding host]
\lean{Turing.FinTM.emLoopHost_input_rewind}
\label{Turing.FinTM.emLoopHost_input_rewind}
\leanok
\uses{Turing.Cfg, Turing.FinTM, Turing.FinTM.emLoopHost, Turing.FinTM.loopControl_idle, Turing.FinTM.loop_rewind_bounded, Turing.MultiTapeTM.runFrom}
\difficulty{3}
Let $c$ be a configuration of the forwarding host in controller phase four. Then there
is a time $t$ at most the input-head position of $c$ plus two after which the host sits
in the startup body-call state with input head at position one and every other component
of $c$ unchanged.
\end{lemma}

\begin{lemma}[Fuel rewind phase in the forwarding host]
\lean{Turing.FinTM.emLoopHost_fuel_rewind}
\label{Turing.FinTM.emLoopHost_fuel_rewind}
\leanok
\uses{Turing.Action.apply, Turing.Cfg, Turing.Cfg.workTapeSymbols, Turing.FinTM, Turing.FinTM.LoopHostState, Turing.FinTM.bufferTape, Turing.FinTM.bufferTape_left, Turing.FinTM.bufferTape_nat, Turing.FinTM.emLoopHost, Turing.FinTM.loopControlAction, Turing.FinTM.loopControl_apply, Turing.FinTM.loopFrame, Turing.FinTM.loopFrame_payload, Turing.FinTM.loopWrite, Turing.MultiTapeTM.runFrom, Turing.MultiTapeTM.runFrom_succ_eq_step, Turing.MultiTapeTM.runFrom_zero, Turing.MultiTapeTM.step}
\difficulty{5}
Let $w$ be a stored word and $j \le |w|$. Starting the forwarding host in its fuel-
rewind phase on a controller frame whose payload track is the buffer tape of $w$ with
payload head at cell $j - 1$, after exactly $j + 1$ steps the host is in the counter-
copy phase with the payload head at cell zero and every other frame component unchanged.
\end{lemma}

\begin{lemma}[Fuel copy phase in the forwarding host]
\lean{Turing.FinTM.emLoopHost_fuel_copy}
\label{Turing.FinTM.emLoopHost_fuel_copy}
\leanok
\uses{Turing.Action.apply, Turing.Cfg, Turing.Cfg.workTapeSymbols, Turing.FinTM, Turing.FinTM.LoopHostState, Turing.FinTM.bufferTape, Turing.FinTM.bufferTape_append, Turing.FinTM.emLoopHost, Turing.FinTM.loopControlAction, Turing.FinTM.loopControl_apply, Turing.FinTM.loopCopyTape, Turing.FinTM.loopCopy_erase, Turing.FinTM.loopCopy_final, Turing.FinTM.loopCopy_read, Turing.FinTM.loopFrame, Turing.FinTM.loopFrame_payload, Turing.FinTM.loopWrite, Turing.MultiTapeTM.runFrom, Turing.MultiTapeTM.runFrom_succ_eq_step, Turing.MultiTapeTM.runFrom_zero, Turing.MultiTapeTM.step}
\difficulty{5}
Let $p$ be the already-copied prefix and $r$ the remaining fuel suffix. Starting the
forwarding host in its counter-copy phase on a controller frame whose counter track is
the buffer tape of $p$ and whose payload track is the partially cleared copying tape,
with both heads at cell $|p|$, after exactly $|r| + 1$ steps the host is in the
synchronized-rewind phase with the counter track holding the full fuel word, the payload
track blank, and both heads one cell left of the word's right end.
\end{lemma}

\begin{lemma}[Synchronized rewind in the forwarding host]
\lean{Turing.FinTM.emLoopHost_fuel_return}
\label{Turing.FinTM.emLoopHost_fuel_return}
\leanok
\uses{Turing.Action.apply, Turing.Cfg, Turing.Cfg.workTapeSymbols, Turing.FinTM, Turing.FinTM.LoopHostState, Turing.FinTM.bufferTape, Turing.FinTM.bufferTape_left, Turing.FinTM.bufferTape_nat, Turing.FinTM.emLoopHost, Turing.FinTM.loopControlAction, Turing.FinTM.loopControl_apply, Turing.FinTM.loopFrame, Turing.FinTM.loopFrame_counter, Turing.FinTM.loopWrite, Turing.MultiTapeTM.runFrom, Turing.MultiTapeTM.runFrom_succ_eq_step, Turing.MultiTapeTM.runFrom_zero, Turing.MultiTapeTM.step}
\difficulty{5}
Let $w$ be the copied counter word and $j \le |w|$. Starting the forwarding host in its
synchronized-rewind phase on a controller frame whose counter track is the buffer tape
of $w$ and whose payload track is blank, with both heads at cell $j - 1$, after exactly
$j + 1$ steps the host is in the input-rewind phase with both heads at cell zero and
every other frame component unchanged.
\end{lemma}

\begin{lemma}[Fuel setup of the forwarding host]
\lean{Turing.FinTM.emLoopHost_fuel_setup}
\label{Turing.FinTM.emLoopHost_fuel_setup}
\leanok
\uses{Turing.Action.apply, Turing.Cfg, Turing.FinTM, Turing.FinTM.LoopHostState, Turing.FinTM.bufferTape, Turing.FinTM.emLoopHost, Turing.FinTM.emLoopHost_fuel_copy, Turing.FinTM.emLoopHost_fuel_return, Turing.FinTM.emLoopHost_fuel_rewind, Turing.FinTM.loopControlAction, Turing.FinTM.loopControl_apply, Turing.FinTM.loopCopy_initial, Turing.FinTM.loopFrame, Turing.FinTM.loopWrite, Turing.MultiTapeTM.runFrom, Turing.MultiTapeTM.runFrom_add, Turing.MultiTapeTM.runFrom_succ_eq_step, Turing.MultiTapeTM.step}
\difficulty{4}
Let $w$ be the captured fuel word. Starting the forwarding host in controller phase zero
on a frame whose counter track is blank and whose payload track is the buffer tape of
$w$ with payload head at cell $|w|$, after exactly $3|w| + 4$ steps the host is in the
input-rewind phase with the counter track holding $w$, the payload track blank, both
heads at cell zero, and all inactive data unchanged.
\end{lemma}

\begin{lemma}[Preparation of the forwarding host]
\lean{Turing.FinTM.emLoopHost_prepare}
\label{Turing.FinTM.emLoopHost_prepare}
\leanok
\uses{Turing.Cfg, Turing.Cfg.Halted, Turing.Cfg.init, Turing.FinTM, Turing.FinTM.ComputesFunInTime, Turing.FinTM.bufferTape, Turing.FinTM.emLoopHost, Turing.FinTM.emLoopHost_fuel_capture, Turing.FinTM.emLoopHost_fuel_setup, Turing.FinTM.emLoopHost_init, Turing.FinTM.emLoopHost_input_rewind, Turing.FinTM.loopFrame, Turing.FinTM.loopFuelCaptured, Turing.FinTM.loopFuelCaptured_frame, Turing.FinTM.loopFuelCfg, Turing.FinTM.loopFuel_init, Turing.FinTM.loopFuel_run, Turing.FinTM.loopReady, Turing.FinTM.loop_first_halt, Turing.FinTM.loop_fuel_width, Turing.FinTM.loop_input_run_le, Turing.FinTM.rightCfg, Turing.MultiTapeTM.initCfg, Turing.MultiTapeTM.runFrom, Turing.MultiTapeTM.runFrom_add}
\difficulty{5}
Let the fuel machine compute, in time $T$, the function sending each input $x$ to the
binary digit word of $R(|x|)$. Then for every input $x$ there are a halted fuel
configuration $c$ with output exactly that digit word and a time $t \le 5\,T(|x|) + 7$
such that the forwarding host, run for $t$ steps from its initial configuration on $x$,
reaches the prepared startup configuration built from $c$.
\end{lemma}

\begin{lemma}[Startup release in the forwarding host]
\lean{Turing.FinTM.emLoopHost_release}
\label{Turing.FinTM.emLoopHost_release}
\leanok
\uses{Turing.Action.apply, Turing.Cfg, Turing.FinTM, Turing.FinTM.emLoopHost, Turing.FinTM.leftCfg, Turing.FinTM.loopBodyCfg, Turing.FinTM.loopBodyPadded, Turing.FinTM.loopCall, Turing.FinTM.loopControlAction, Turing.MultiTapeTM.step, Turing.captureCfg, Turing.moveInputPos_zero}
\difficulty{4}
Let $c$ be a body configuration, $w$ a counter word, and $f$ a fuel configuration. One
step of the forwarding host from the startup call on $c$ with control removed, release
clear, and flag $\mathrm{false}$ yields the active call on $c$ with control the anchor,
release set, and no flag, changing no body, counter, or fuel data.
\end{lemma}

\begin{lemma}[One borrow step in the forwarding host]
\lean{Turing.FinTM.emLoopHost_borrow_step}
\label{Turing.FinTM.emLoopHost_borrow_step}
\leanok
\uses{Turing.Action.apply, Turing.Cfg, Turing.Cfg.workTapeSymbols, Turing.FinTM, Turing.FinTM.LoopHostState, Turing.FinTM.bufferTape, Turing.FinTM.emLoopHost, Turing.FinTM.loopBuffer_read, Turing.FinTM.loopBuffer_write, Turing.FinTM.loopControlAction, Turing.FinTM.loopControl_apply, Turing.FinTM.loopFrame, Turing.FinTM.loopFrame_counter, Turing.FinTM.loopWrite, Turing.MultiTapeTM.step}
\difficulty{4}
Let $p$ and $r$ be Boolean words. One step of the forwarding host from its borrow phase
on a frame whose counter track is the buffer tape of $p$ followed by $r$ with counter
head at $|p|$ behaves as the standalone decrement machine: an empty $r$ enters the
underflow rewind, a leading $\mathrm{true}$ writes $\mathrm{false}$ and enters the
success rewind, and a leading $\mathrm{false}$ writes $\mathrm{true}$ and continues the
borrow. Only the counter track changes.
\end{lemma}

\begin{lemma}[The borrow scan in the forwarding host]
\lean{Turing.FinTM.emLoopHost_borrow_run}
\label{Turing.FinTM.emLoopHost_borrow_run}
\leanok
\uses{Turing.Cfg, Turing.FinTM, Turing.FinTM.LoopHostState, Turing.FinTM.bufferTape, Turing.FinTM.emLoopHost, Turing.FinTM.emLoopHost_borrow_step, Turing.FinTM.loopBorrowPos, Turing.FinTM.loopDebit, Turing.FinTM.loopFrame, Turing.MultiTapeTM.runFrom, Turing.MultiTapeTM.runFrom_succ_eq_step}
\difficulty{4}
Let $p$ and $w$ be Boolean words and $j$ the length of the maximal $\mathrm{false}$
prefix of $w$. Running the forwarding host for exactly $j + 1$ steps from its borrow
phase on a frame whose counter track holds $p$ followed by $w$ with head at $|p|$
reaches the success- or underflow-rewind phase according to the decrement's flag, with
the counter track holding $p$ followed by the decremented word and head at $|p| + j -
1$.
\end{lemma}

\begin{lemma}[Borrow rewinds in the forwarding host]
\lean{Turing.FinTM.emLoopHost_borrow_rewind}
\label{Turing.FinTM.emLoopHost_borrow_rewind}
\leanok
\uses{Turing.Action.apply, Turing.Cfg, Turing.Cfg.workTapeSymbols, Turing.FinTM, Turing.FinTM.LoopHostState, Turing.FinTM.bufferTape, Turing.FinTM.bufferTape_left, Turing.FinTM.bufferTape_nat, Turing.FinTM.emLoopHost, Turing.FinTM.loopControlAction, Turing.FinTM.loopControl_apply, Turing.FinTM.loopFrame, Turing.FinTM.loopFrame_counter, Turing.FinTM.loopWrite, Turing.MultiTapeTM.runFrom, Turing.MultiTapeTM.runFrom_succ_eq_step, Turing.MultiTapeTM.runFrom_zero, Turing.MultiTapeTM.step}
\difficulty{5}
Let $w$ be a counter word, $b$ a success bit, and $j \le |w|$. Running the forwarding
host for exactly $j + 1$ steps from the rewind phase selected by $b$ on a frame whose
counter track is the buffer tape of $w$ with head at cell $j - 1$ returns the counter
head to zero and dispatches: success releases the next anchor call, underflow enters the
exhaustion phase without yet emitting.
\end{lemma}

\begin{lemma}[Complete counter operation in the forwarding host]
\lean{Turing.FinTM.emLoopHost_borrow}
\label{Turing.FinTM.emLoopHost_borrow}
\leanok
\uses{Turing.Cfg, Turing.FinTM, Turing.FinTM.LoopHostState, Turing.FinTM.bufferTape, Turing.FinTM.emLoopHost, Turing.FinTM.emLoopHost_borrow_rewind, Turing.FinTM.emLoopHost_borrow_run, Turing.FinTM.loopBorrowPos, Turing.FinTM.loopBorrowPos_le, Turing.FinTM.loopDebit, Turing.FinTM.loopDebit_length, Turing.FinTM.loopFrame, Turing.MultiTapeTM.runFrom, Turing.MultiTapeTM.runFrom_add}
\difficulty{4}
Let $w$ be a counter word and $j$ the length of its maximal $\mathrm{false}$ prefix.
Then $2j + 2 \le 2|w| + 2$, and running the forwarding host for exactly $2j + 2$ steps
from its borrow phase on a frame whose counter track is the buffer tape of $w$ with head
at zero ends with the counter track holding the decremented word, head at zero, in the
released next-anchor call state on success and in the exhaustion phase on underflow.
\end{lemma}

\begin{lemma}[A rejecting round in the forwarding host]
\lean{Turing.FinTM.emLoopHost_reject}
\label{Turing.FinTM.emLoopHost_reject}
\leanok
\uses{Turing.Action.apply, Turing.Cfg, Turing.Cfg.workTapeSymbols, Turing.FinTM, Turing.FinTM.bufferTape, Turing.FinTM.emLoopHost, Turing.FinTM.emLoopHost_borrow, Turing.FinTM.leftCfg, Turing.FinTM.loopBodyCfg, Turing.FinTM.loopBodyPadded, Turing.FinTM.loopBorrowPos, Turing.FinTM.loopBorrowPos_le, Turing.FinTM.loopCall, Turing.FinTM.loopCall_frame, Turing.FinTM.loopCall_reframe, Turing.FinTM.loopControlAction, Turing.FinTM.loopControl_apply, Turing.FinTM.loopDebit, Turing.FinTM.loopFlag_clear, Turing.FinTM.loopFrame, Turing.FinTM.loopFrame_flag, Turing.FinTM.loopWrite, Turing.MultiTapeTM.runFrom, Turing.MultiTapeTM.runFrom_succ_eq_step, Turing.MultiTapeTM.runFrom_succ_eq_step', Turing.MultiTapeTM.step, Turing.captureCfg}
\difficulty{5}
Let $c$ be a halted body configuration with empty output, $w$ a counter word, and $f$ a
fuel configuration. Then there is a time $t \le 2|w| + 4$ such that, from the stopped
call on $c$ with stop flag $\mathrm{false}$: if the fixed-width decrement of $w$
succeeds, the forwarding host at time $t$ is at the released next-anchor call with the
decremented counter word; otherwise the host at time $t$ has halted with output empty in
find mode and the single bit $\mathrm{false}$ in decision mode.
\end{lemma}

\begin{definition}[Forwarded body configuration]
\lean{Turing.FinTM.emLoopForwardCfg}
\label{Turing.FinTM.emLoopForwardCfg}
\leanok
\uses{Turing.Cfg, Turing.FinTM.bufferTape, Turing.FinTM.leftCfg, Turing.emitCfg}
For a control embedding, a return state, an accumulated output prefix $p$, and a source
configuration $c$, the host configuration that embeds $c$ with one extra blank inactive
tape, carries $p$ followed by the output of $c$ as physical output, and embeds the
control of $c$, using the return state when $c$ has halted.
\end{definition}

\begin{lemma}[A forwarded action acts on a forwarded frame]
\lean{Turing.FinTM.emLoop_forward_apply}
\label{Turing.FinTM.emLoop_forward_apply}
\leanok
\uses{Turing.Action, Turing.Action.apply, Turing.Cfg, Turing.FinTM.emLoopForwardCfg, Turing.FinTM.leftAction, Turing.FinTM.leftCfg_apply, Turing.emitAction, Turing.emitCfg}
\difficulty{3}
For every control embedding, return state, output prefix $p$, source action $a$, and
source configuration $c$: applying the padded forwarding action of $a$ to the forwarded
configuration of $c$ equals the forwarded configuration of $a$ applied to $c$, appending
the optional emitted bit after the accumulated prefix, including on a halting action.
\end{lemma}

\begin{lemma}[Guarded forwarding with one inactive tape]
\lean{Turing.FinTM.emLoop_forward_run}
\label{Turing.FinTM.emLoop_forward_run}
\leanok
\uses{Turing.Cfg, Turing.Cfg.inputSymbol, Turing.Cfg.workTapeSymbols, Turing.FinTM.emLoopForwardCfg, Turing.FinTM.emLoop_forward_apply, Turing.FinTM.leftAction, Turing.FinTM.leftCfg, Turing.MultiTapeTM, Turing.MultiTapeTM.runFrom, Turing.MultiTapeTM.runFrom_succ_eq_step', Turing.MultiTapeTM.step, Turing.emitAction, Turing.emitCfg}
\difficulty{4}
Let a source machine and a host machine be given, together with a control embedding and
a return state, such that on every embedded state the host transition is the padded
forwarding of the source transition. If the source run from a configuration $c$ stays
live before time $t$, then the host run for $t$ steps from the forwarded configuration
of $c$ with prefix $p$ equals the forwarded configuration of the source run at time $t$
with the same prefix; the source may halt on the final action, whose emission is
forwarded.
\end{lemma}

\begin{definition}[A forwarding body call]
\lean{Turing.FinTM.emLoopCall}
\label{Turing.FinTM.emLoopCall}
\leanok
\uses{Turing.Cfg, Turing.FinTM, Turing.FinTM.LoopHostState, Turing.FinTM.emLoopForwardCfg, Turing.FinTM.loopBodyPadded}
For an output prefix $p$, a startup marker, a body configuration $c$, a release bit, an
optional flag bit, a counter word, and a fuel configuration, the forwarding-host
configuration whose body, counter, and fuel data are those of the padded stopped-body
configuration under the forwarding discipline, whose physical output is $p$ followed by
the output of $c$, and whose now-unused payload tape is blank with head at the origin.
\end{definition}

\begin{lemma}[A forwarding call as a clean call plus prefix]
\lean{Turing.FinTM.emLoopCall_frame}
\label{Turing.FinTM.emLoopCall_frame}
\leanok
\uses{Turing.Cfg, Turing.FinTM, Turing.FinTM.emLoopCall, Turing.FinTM.emLoopForwardCfg, Turing.FinTM.leftCfg, Turing.FinTM.loopBodyCfg, Turing.FinTM.loopBodyPadded, Turing.FinTM.loopCall, Turing.captureCfg, Turing.emitCfg}
\difficulty{4}
For every output prefix $p$, startup marker, body configuration $c$, release bit, flag
bit, counter word, and fuel configuration: the forwarding call equals the corresponding
capturing call built on $c$ with its output removed, with physical output replaced by
$p$ followed by the output of $c$. The identity is between configurations only.
\end{lemma}

\begin{lemma}[Forwarding of the stopped-body run]
\lean{Turing.FinTM.emLoopHost_body_forward}
\label{Turing.FinTM.emLoopHost_body_forward}
\leanok
\uses{Turing.Cfg, Turing.Cfg.Halted, Turing.FinTM, Turing.FinTM.emLoopForwardCfg, Turing.FinTM.emLoopHost, Turing.FinTM.emLoop_forward_run, Turing.FinTM.loopBodySource, Turing.MultiTapeTM.runFrom}
\difficulty{1}
Let $c$ be a configuration of the extended stopped-body source on the forwarding host's
layout, $p$ an output prefix, and $t$ a time before which that source does not halt.
Then the forwarding host run for $t$ steps from the forwarded embedding of $c$ with
prefix $p$ equals the forwarded embedding of the source run at time $t$ with the same
prefix.
\end{lemma}

\begin{lemma}[Anchor returns in the forwarding host]
\lean{Turing.FinTM.emLoopHost_anchor_return}
\label{Turing.FinTM.emLoopHost_anchor_return}
\leanok
\uses{Turing.Cfg, Turing.Cfg.Halted, Turing.FinTM, Turing.FinTM.emLoopCall, Turing.FinTM.emLoopHost, Turing.FinTM.emLoopHost_body_forward, Turing.FinTM.leftCfg, Turing.FinTM.loopBodyCfg, Turing.FinTM.loopBodyPadded, Turing.FinTM.loopBodySource_run, Turing.FinTM.loopBodyTM, Turing.FinTM.loopBody_run, Turing.FinTM.loopBody_stop, Turing.FinTM.loop_live_prefix, Turing.MultiTapeTM.runFrom, Turing.MultiTapeTM.runFrom_succ_eq_step'}
\difficulty{5}
Let $c$ be a body configuration, $p$ an output prefix, and $t$ a time at which the body
run from $c$ sits in the anchor state, with no earlier anchor visit except possibly a
released step at time zero and with the release bit clear whenever $t = 0$. Then the
forwarding host, run for $t + 1$ steps from the forwarding call on $c$ with prefix $p$,
reaches the forwarding stopped call on the body run at time $t$ with control removed,
release clear, and stop flag $\mathrm{false}$, the counter word and fuel residue
unchanged.
\end{lemma}

\begin{lemma}[Empty-prefix forwarding calls are capturing calls]
\lean{Turing.FinTM.emLoopCall_empty}
\label{Turing.FinTM.emLoopCall_empty}
\leanok
\uses{Turing.Cfg, Turing.FinTM, Turing.FinTM.emLoopCall, Turing.FinTM.emLoopCall_frame, Turing.FinTM.loopCall}
\difficulty{2}
For every startup marker, body configuration $c$ with empty output, release bit, flag
bit, counter word, and fuel configuration, the forwarding call on $c$ with empty prefix
equals the corresponding capturing call on $c$.
\end{lemma}

\begin{lemma}[Silent startup of the forwarding host]
\lean{Turing.FinTM.emLoopHost_start}
\label{Turing.FinTM.emLoopHost_start}
\leanok
\uses{Turing.Cfg, Turing.Cfg.ofWords, Turing.FinTM, Turing.FinTM.emLoopCall_empty, Turing.FinTM.emLoopHost, Turing.FinTM.emLoopHost_anchor_return, Turing.FinTM.emLoopHost_release, Turing.FinTM.loopCall, Turing.FinTM.loopReady, Turing.FinTM.loopReady_call, Turing.MultiTapeTM.initCfg, Turing.MultiTapeTM.runFrom, Turing.MultiTapeTM.runFrom_succ_eq_step', Turing.stateWord}
\difficulty{3}
Let $s$ be a word, $f$ a fuel configuration, and $t$ a time such that the body, run from
its initial configuration, first visits the anchor at time $t$, arriving at the clean
seam that carries $s$ on work tape zero with all other work tapes blank. Then the
forwarding host, run for $t + 2$ steps from the prepared startup configuration, reaches
the released call on that seam with counter word the output of $f$ and fuel residue $f$,
with nothing emitted and no fuel debited.
\end{lemma}

\begin{lemma}[One emitting round of the forwarding host]
\lean{Turing.FinTM.emLoopHost_round}
\label{Turing.FinTM.emLoopHost_round}
\leanok
\uses{Turing.Cfg, Turing.Cfg.ofWords, Turing.FinTM, Turing.FinTM.emLoopCall_empty, Turing.FinTM.emLoopCall_frame, Turing.FinTM.emLoopHost, Turing.FinTM.emLoopHost_anchor_return, Turing.FinTM.emLoopHost_reject, Turing.FinTM.emLoop_run_prefix, Turing.FinTM.loopCall, Turing.FinTM.loopDebit, Turing.MultiTapeTM.runFrom, Turing.MultiTapeTM.runFrom_add, Turing.stateWord}
\difficulty{5}
Let $s$ be a round word, $w$ a counter word, $f$ a fuel configuration, and $t > 0$ a
time such that the body, started at the clean anchor seam carrying $s$, does not revisit
the anchor strictly between times zero and $t$ and at time $t$ sits at the clean seam
carrying the stepped word with output a chunk word $u$. Then there is a time $v \le t +
2|w| + 5$ such that from the released call on the seam of $s$: when the fixed-width
decrement of $w$ succeeds, the forwarding host at time $v$ is at the released call on
the next seam with the decremented counter and physical output exactly $u$; on underflow
the host at time $v$ has halted with output exactly $u$ and no verdict bit.
\end{lemma}

\begin{lemma}[Prefix summation of emitting segments]
\lean{Turing.FinTM.emLoop_sum}
\label{Turing.FinTM.emLoop_sum}
\leanok
\uses{Turing.Cfg, Turing.FinTM.emLoop_run_prefix, Turing.MultiTapeTM, Turing.MultiTapeTM.runFrom, Turing.MultiTapeTM.runFrom_add}
\difficulty{6}
Let $M$ be a multi-tape Turing machine over the Boolean alphabet, let $c_0, \dots,
c_{N-1}$ be configurations on a common input with empty outputs, let $u_0, \dots,
u_{N-1}$ be chunk words, and suppose $N > 0$ and that for each $i < N$, within $B$
steps, the machine run from $c_i$ reaches $c_{i+1}$ with output exactly $u_i$ when $i +
1 < N$, and halts with output exactly $u_i$ when $i + 1 = N$. Then for every accumulated
prefix $p$, the machine run from $c_0$ with output $p$ halts within $N \cdot B$ steps
with output $p$ followed by the concatenation $u_0 u_1 \cdots u_{N-1}$ in order.
\end{lemma}

\begin{theorem}[The emitting loop combinator]
\lean{Turing.FinTM.exists_emitLoopTM}
\label{Turing.FinTM.exists_emitLoopTM}
\leanok
\uses{Turing.Cfg.ofWords, Turing.FinTM, Turing.FinTM.ComputesFunInTime, Turing.FinTM.ComputesInTime, Turing.FinTM.ComputesInTime.mono, Turing.FinTM.emLoopHost, Turing.FinTM.emLoopHost_prepare, Turing.FinTM.emLoopHost_round, Turing.FinTM.emLoopHost_start, Turing.FinTM.emLoop_sum, Turing.FinTM.loopCall, Turing.FinTM.loopDebit, Turing.FinTM.loopDebit_iterate_length, Turing.FinTM.loopDebit_iterate_value, Turing.FinTM.loopDebit_success, Turing.FinTM.loop_fuel_width, Turing.FinTM.loop_orbit_inv, Turing.MultiTapeTM.initCfg, Turing.MultiTapeTM.runFrom, Turing.MultiTapeTM.runFrom_add, Turing.stateWord}
\difficulty{6}
Let a loop body with anchor be given, a fuel machine computing the binary digits of
$R(|x|)$ in time $T$, an invariant that holds at the initial word and is preserved by
the step function, a startup reaching the initial seam within $T(|x|)$ without visiting
the anchor earlier, and a per-round contract: on every admissible word $s$, in positive
time at most $T(|x|)$ with no early anchor revisit, the body moves from the clean seam
of $s$ to the clean seam of the stepped word while appending exactly the chunk word
assigned to $(x, s)$ to its output. Then there are a finite machine $E$ and a constant
$c$ such that $E$ computes, within $c (T(n) + 1)(R(n) + 2)$ steps on inputs of length
$n$, the concatenation of the chunks of the first $R(|x|) + 1$ orbit points of the
initial word, in order.
\end{theorem}

\begin{theorem}[The install-mode clean call]
\lean{Turing.FinTM.exists_installCallTM}
\label{Turing.FinTM.exists_installCallTM}
\leanok
\uses{Turing.Cfg.ofWords, Turing.FinTM, Turing.FinTM.ComputesFunInTime, Turing.FinTM.emCallTM, Turing.FinTM.emCall_first, Turing.MultiTapeTM.runFrom, Turing.stateWord}
\difficulty{2}
Let $M$ be a bundled finite machine computing a function $f$ on Boolean words in time
$T$. Then there are a machine $C$ with a positive number of work tapes, control states
called entry and exit, and a constant $c$ such that for every native input $x$ and
argument $a$: starting $C$ at the canonical seam in the entry state with $a$ as the sole
tape-resident word (all other tapes blank, input head at one, empty output), there is a
positive time $t \le c\,(T(|a|) + |a| + |f(a)| + 1)$, before which the exit state is
never visited, at which $C$ is exactly at the canonical seam in the exit state carrying
$f(a)$ as the sole tape-resident word: the result installed, all scratch restored,
nothing emitted, and the native input untouched.
\end{theorem}

\begin{theorem}[The emit-mode clean call]
\lean{Turing.FinTM.exists_emitCallTM}
\label{Turing.FinTM.exists_emitCallTM}
\leanok
\uses{Turing.Cfg.ofWords, Turing.FinTM, Turing.FinTM.ComputesFunInTime, Turing.FinTM.emCallTM, Turing.FinTM.emCall_first, Turing.MultiTapeTM.runFrom, Turing.stateWord}
\difficulty{2}
Let $M$ be a bundled finite machine computing a function $f$ on Boolean words in time
$T$. Then there are a machine $C$ with a positive number of work tapes, control states
called entry and exit, and a constant $c$ such that for every native input $x$ and
argument $a$: starting $C$ at the canonical seam in the entry state with $a$ as the sole
tape-resident word, there is a positive time $t \le c\,(T(|a|) + |a| + |f(a)| + 1)$,
before which the exit state is never visited, at which $C$ is exactly at the entry seam
with exit control and physical output $f(a)$: the argument preserved on its tape, all
scratch restored, and the computed word emitted.
\end{theorem}
